\documentclass[12pt,letterpaper]{article}
\usepackage[T1]{fontenc}
\usepackage{mathptmx}
\usepackage{courier}
\usepackage[letterpaper,margin=1in]{geometry}
\usepackage{amsmath,amssymb,amsthm}
\usepackage{microtype}
\usepackage{booktabs}
\usepackage{array}
\usepackage{enumitem}
\usepackage[hidelinks]{hyperref}

\numberwithin{equation}{section}

\theoremstyle{plain}
\newtheorem{theorem}{Theorem}[section]
\newtheorem{proposition}[theorem]{Proposition}
\newtheorem{lemma}[theorem]{Lemma}
\newtheorem{corollary}[theorem]{Corollary}
\theoremstyle{definition}
\newtheorem{definition}[theorem]{Definition}
\newtheorem{example}[theorem]{Example}
\newtheorem{counterexample}[theorem]{Counterexample}
\newtheorem{remark}[theorem]{Remark}
\newtheorem{question}[theorem]{Question}
\newtheorem{yremark}{Remark}
\renewcommand{\theyremark}{2026.\arabic{yremark}}

\newcommand{\F}{\mathbb{F}}
\newcommand{\Z}{\mathbb{Z}}
\newcommand{\C}{\mathbb{C}}
\newcommand{\Q}{\mathbb{Q}}
\DeclareMathOperator{\Cent}{Cent}
\DeclareMathOperator{\Conj}{Conj}
\DeclareMathOperator{\Ind}{Ind}
\DeclareMathOperator{\Res}{Res}
\DeclareMathOperator{\supp}{supp}
\DeclareMathOperator{\GL}{GL}
\DeclareMathOperator{\SL}{SL}
\DeclareMathOperator{\PGL}{PGL}
\DeclareMathOperator{\PSL}{PSL}
\DeclareMathOperator{\AGL}{AGL}
\DeclareMathOperator{\diag}{diag}
\DeclareMathOperator{\Span}{Span}
\newcommand{\triv}{\mathbf{1}}

\title{The representation theory of finite commutation transitive groups\\[4pt]
\large Draft 1. Corrected and expanded version (an expository note)}
\author{David Joyner\thanks{This note corrects and expands the unpublished
manuscript ``The representation theory of finite commutation transitive
groups, preliminary version'' (file \texttt{comm\_tr6.tex}, dated 8-5-2000).
Revision prepared November 16, 2026.}\\
\texttt{wdjoyner@gmail.com}}
\date{Original draft: 2000. Corrected and expanded draft: 2026}

\begin{document}
\maketitle

\begin{abstract}
This expository note was revised, corrected and expanded with the help of
an AI system (Claude, by Anthropic). See the Acknowledgments for details.
A finite group is commutative transitive (CT) if commuting is a transitive
relation on its non-identity elements. It is commutation transitive if the
same holds on its non-central elements. Equivalently, every centralizer of a
non-identity (respectively non-central) element is abelian. We give detailed
definitions and examples, and we survey the classification of finite CT
groups. We then study the representations obtained by inducing characters
from centralizers. We prove a corrected induced-character formula and an
exact irreducibility criterion. We also prove that, in a finite CT group,
characters induced from centralizers span the full character ring over
$\Z$. Finally, we connect the subject to graph theory. The map sending an
element to its conjugacy class is a harmonic morphism, in the sense of Baker
and Norine, from the commuting graph onto the commuting conjugacy class
graph. Several claims of the 2000 draft were false. We keep them as worked
counterexamples. One example is the claim that $\PGL(2,\F_q)$, $q$ odd, is
CT.
\end{abstract}

\tableofcontents

%======================================================================
\section{Introduction}\label{sec:intro}

This section explains what the note is about and what it proves.

\subsection{The problem}

In an abelian group every pair of elements commutes. Commutative
transitive groups are the next simplest case. Two non-identity elements may
fail to commute. But if $x$ commutes with $y$, and $y$ commutes with $z$,
then $x$ commutes with $z$. So the non-identity elements fall into
``commuting families'', and each family together with the identity is an
abelian subgroup.

Such groups appear in two quite different parts of group theory.
\begin{itemize}[itemsep=1pt]
\item In finite group theory they are the \emph{CA-groups}. Suzuki's work on
  them in 1957 was a first step toward the odd order theorem
  \cite{Su57,BSW58}.
\item In combinatorial group theory, free groups and many related groups are
  CT. The property plays a role in the study of residually free groups and
  the Tarski problems \cite{Fi83,FR09,FGRS15}.
\end{itemize}
The finite CT groups are completely classified (Theorem~\ref{thm:classif}).
Each one is abelian, a Frobenius group with abelian kernel and cyclic
complement, or $\PSL(2,\F_{2^f})$ with $f\ge 2$.

\subsection{What this note does}

The 2000 draft asked how the representations of a CT group arise from its
centralizers. This note answers that question with complete proofs. The
main results are the following.
\begin{enumerate}[itemsep=1pt]
\item A correct formula for the character of $\Ind_C^G\mu$, where $C$ is a
  centralizer and $\mu$ is a character of $C$
  (Proposition~\ref{prop:indformula}).
\item An exact criterion: in a non-abelian finite CT group, $\Ind_C^G\mu$ is
  irreducible if and only if $C$ is normal and $\mu\neq\triv$
  (Theorem~\ref{thm:irred}).
\item Linear independence of the permutation characters $\triv_A^G$, with
  one $A$ from each conjugacy class of centralizers
  (Proposition~\ref{prop:indep}). We also give the exact dimension of the
  span of the ``support vectors'' studied in 2000
  (Proposition~\ref{prop:taudim}).
\item A spanning theorem. In a finite CT group, every virtual character is a
  $\Z$-linear combination of characters induced from linear characters of
  centralizers (Theorem~\ref{thm:span}). The proof uses Brauer's induction
  theorem.
\item A connection with graph theory. The commuting graph of a commutation
  transitive group is a disjoint union of complete graphs. The class map from
  it onto the commuting conjugacy class graph is a harmonic morphism
  (Theorem~\ref{thm:harmonic}). On the non-commuting graph, the analogous
  map is harmonic only when all centralizers of non-central elements have
  the same order (Proposition~\ref{prop:noncommharm}).
\end{enumerate}
Several statements of the 2000 draft were wrong. We do not hide them. Each
one is stated and refuted as a counterexample, because each one teaches
something. They are Counterexamples~\ref{cex:pgl}, \ref{cex:quotient},
\ref{cex:formula}, \ref{cex:pglcor}, \ref{cex:supports} and~\ref{cex:q8}.

\subsection{Scope and conventions}

All groups are finite, except in the literature review and in
Question~\ref{q:profinite}. The 2000 draft announced a profinite analogue,
but stated and proved nothing. That section has been removed. All
representations are complex and finite-dimensional. We identify a
representation with its character when no confusion can arise.

Each result is labeled by its status. A statement with a proof is proved
here. A statement marked ``cited'' is taken from the given source. Numbered
remarks of the form ``Remark~2026.$n$'' contain material that we have not
fully verified. Some claims were also checked by computer. The script
\texttt{comm\_tr\_checks.py} accompanies this note. It is in the folder
\texttt{commutation-transitive/verification} of the repository
\url{https://github.com/wdjoyner/alg-graph-thry}. It is exhaustive for the
small groups it treats, and the note says when a claim depends on it.

%======================================================================
\section{Literature review}\label{sec:lit}

This section summarizes the literature and says where the present note fits.
We separate work before and after the original draft of 2000.

\subsection{Before 2000}

\paragraph{Weisner (1925).} Weisner \cite{We25} studied groups ``in which
the normaliser of every element except identity is abelian''. In modern
language, ``normaliser of an element'' means its centralizer. So these are
exactly the CT groups. He showed that a finite CT group is solvable or
simple. Wu later observed that one of Weisner's arguments has a gap
\cite{Wu98}.

\paragraph{Suzuki (1957) and Brauer--Suzuki--Wall (1958).} Suzuki
\cite{Su57} proved that a CA-group of odd order is solvable. This was the
first serious attack on Burnside's conjecture, later proved as the
Feit--Thompson theorem. His argument is character-theoretic. It applies the
orthogonality relations to the abelian centralizers and their
normalizers. Combined with the work of Brauer, Suzuki and Wall
\cite{BSW58}, it shows that the groups $\PSL(2,\F_{2^n})$ are the only
non-abelian simple groups with this property. A textbook account of the
classification is in Suzuki's \emph{Group Theory II} \cite{Su86}.

\paragraph{Schmidt (1970).} Schmidt \cite{Sc70} studied finite groups whose
lattice of centralizers is distributive or modular. Later authors use his
results to describe the finite groups in which centralizers of
\emph{non-central} elements are abelian
\cite{GM23}. These are our commutation transitive groups, which those
authors call \emph{AC-groups}.

\paragraph{Fine (1983) and Levin--Rosenberger (1986).} The name
``commutative transitive'' was coined by Fine \cite{Fi83}, in the setting of
Fuchsian and Kleinian groups. This attribution follows \cite{FGRS15}. Levin
and Rosenberger \cite{LR86} studied power-commutative and
commutation-transitive groups.

\paragraph{Wu (1998).} Wu \cite{Wu98} gave a correct proof of Weisner's
solvable case. A finite solvable CT group is the semidirect product of its
Fitting subgroup and a cyclic group acting fixed-point-freely. Wu also
classified the locally finite CT groups, and noted that Tarski monsters are
infinite simple CT groups.

\paragraph{Representation theory.} The tools used here are classical. They
are induced representations, Frobenius reciprocity, Mackey's criteria, and
Brauer's induction theorem. They are explained in Serre's book \cite{Se77}.
For Frobenius groups we use Isaacs \cite{Is76}.

\subsection{After 2000}

\paragraph{The original draft.} The preliminary version of this note was
circulated in 2000. A companion preprint computed examples for $A_5$ in GAP
\cite{Jo00}.

\paragraph{CT groups in combinatorial and model-theoretic group theory.}
Fine, Gaglione, Rosenberger and Spellman studied the ``commutative
transitive kernel'' \cite{FGRS00}. Fine and Rosenberger surveyed the subject
in 2009 \cite{FR09}. In 2015, Fine, Gaglione, Rosenberger and Spellman
compared CT groups with CSA groups, in which maximal abelian subgroups are
malnormal \cite{FGRS15}. They explain that commutative transitivity was
important in the solution of the Tarski problems. Duncan, Kazachkov and
Remeslennikov showed that the CT groups are exactly the groups with
trivial center and centralizer dimension~$2$ \cite{DKR05}.

\paragraph{Commuting and non-commuting graphs.} Abdollahi, Akbari and
Maimani introduced the non-commuting graph of a group, on the vertex set
$G\setminus Z(G)$ \cite{AAM06}. AC-groups play a special role there. Grazian
and Monetta continued this study \cite{GM23}. The commuting graph of
Section~\ref{sec:graph} is the complement of the non-commuting graph.

\paragraph{Conjugacy class graphs.} Herzog, Longobardi and Maj introduced
the commuting conjugacy class graph in 2009 \cite{HLM09}. Its vertices are
the non-central conjugacy classes, and two classes are adjacent when they
contain commuting elements. Rezaei and Foruzanfar determined this graph when
$G/Z(G)$ is a Frobenius group of order $pq$ or $p^2q$ \cite{RF23}. For
commutation transitive groups, Section~\ref{sec:harmonic} shows that it is a
disjoint union of complete graphs.

\paragraph{Harmonic morphisms of graphs.} Urakawa introduced harmonic
morphisms between simple graphs \cite{Ur00}. Baker and Norine adapted the
definition to multigraphs. They developed the theory as a discrete analogue
of holomorphic maps between Riemann surfaces, including a Riemann--Hurwitz
formula \cite{BN09}. Corry defined harmonic group actions on graphs and
proved genus bounds for them \cite{Co10}. Melles and Joyner explore graphs
through harmonic morphisms in two books \cite{JM17,MJ26}.
Section~\ref{sec:harmonic} applies this theory to the commuting graph.

\paragraph{Rings.} Dol\v{z}an, Klep and Moravec studied finite rings in which
commutativity is transitive \cite{DKM09}. Their introduction credits Wu
\cite{Wu98} with proving that finite CT groups are solvable or simple.

\paragraph{How this note fits.} We found no published treatment of the
specific questions asked here. These are which characters induced from
centralizers are irreducible, how they span the character ring, and whether
the class map on the commuting graph is harmonic. All
results in Sections~\ref{sec:induced}--\ref{sec:harmonic} are elementary
consequences of the classification and standard tools. Their value is
expository. They show how the abelian centralizers organize the
representation theory.

%======================================================================
\section{Background}\label{sec:background}

This section collects the group theory and representation theory used
later. Each subsection has an example.

\subsection{Centralizers, normalizers and conjugacy}\label{sec:conj}

Let $G$ be a finite group. The \emph{centralizer} of $g\in G$ is
$\Cent(g,G)=\{x\in G : xg=gx\}$. The \emph{center} is $Z=Z(G)=\{g : gx=xg
\text{ for all }x\in G\}$. The \emph{normalizer} of a subgroup $H$ is
$N_G(H)=\{x\in G: xHx^{-1}=H\}$. The \emph{conjugacy class} of $g$ is
$\Conj(g,G)=\{x^{-1}gx : x\in G\}$. We write $G_*$ for the set of conjugacy
classes.

\begin{lemma}\label{lem:conjcent}
For $g,y\in G$ we have $y^{-1}\Cent(g,G)y=\Cent(y^{-1}gy,G)$. The map
$\Cent(g,G)x\mapsto x^{-1}gx$ is a well-defined bijection from the set of
right cosets of $\Cent(g,G)$ onto $\Conj(g,G)$. In particular
$|\Conj(g,G)|=[G:\Cent(g,G)]$.
\end{lemma}

\begin{proof}
If $xg=gx$ then $(y^{-1}xy)(y^{-1}gy)=y^{-1}xgy=y^{-1}gxy=
(y^{-1}gy)(y^{-1}xy)$. This shows one inclusion, and the reverse inclusion
follows by replacing $y$ with $y^{-1}$. For the second claim,
$x^{-1}gx=y^{-1}gy$ holds if and only if $yx^{-1}$ commutes with $g$. That
is, $y\in\Cent(g,G)x$.
\end{proof}

\begin{example}\label{ex:s3conj}
Let $G=S_3$ and $g=(123)$. Then $\Cent(g,G)=A_3$ and $\Conj(g,G)=
\{(123),(132)\}$, which has $2=[S_3:A_3]$ elements. For $g=(12)$ we get
$\Cent(g,G)=\{1,(12)\}$ and three conjugates.

The cosets in Lemma~\ref{lem:conjcent} must be \emph{right} cosets. The 2000
draft used left cosets $x\Cent(g,G)$, and that map is not well defined. Take
$g=(12)$, $x=(13)$ and $c=(12)\in\Cent(g,G)$. Then $x$ and $xc$ lie in the
same left coset. But $x^{-1}gx=(23)$, while $(xc)^{-1}g(xc)=
c^{-1}(23)c=(13)$.
\end{example}

\subsection{Frobenius groups and fixed-point-free actions}\label{sec:frob}

Frobenius groups are the solvable non-abelian CT groups, so we recall them
here.

\begin{definition}\label{def:frob}
A \emph{Frobenius group} is a finite group $G$ with a subgroup $H$ such
that $1\neq H\neq G$ and $H\cap xHx^{-1}=1$ for all $x\in G\setminus H$.
The subgroup $H$ is a \emph{Frobenius complement}. A subgroup with the
intersection property is called \emph{malnormal}.
\end{definition}

\begin{theorem}[Frobenius; cited from {\cite[Thm.~7.2]{Is76}}]
\label{thm:frobenius}
Let $G$ be a Frobenius group with complement $H$. Then
$K=\bigl(G\setminus\bigcup_{x\in G}xHx^{-1}\bigr)\cup\{1\}$ is a normal
subgroup of $G$, called the \emph{Frobenius kernel}. Moreover $G=KH$ and
$K\cap H=1$.
\end{theorem}

An automorphism $\alpha$ of a group $N$ is \emph{fixed-point-free} if
$\alpha(n)=n$ only for $n=1$. A group $H$ acts fixed-point-freely on $N$ if
every non-identity element of $H$ acts by a fixed-point-free automorphism.

\begin{example}\label{ex:frob}
The following are Frobenius groups.
\begin{enumerate}[itemsep=1pt]
\item $S_3$, with kernel $A_3$ and complement $\langle(12)\rangle$.
\item $A_4$, with kernel $V_4=\{1,(12)(34),(13)(24),(14)(23)\}$ and
  complement $\langle(123)\rangle$.
\item The dihedral group $D_{2n}$ of order $2n$ with $n$ odd. The kernel is
  the rotation subgroup and the complement is generated by a reflection.
\item The affine group $\AGL(1,q)=\{x\mapsto ax+b : a\in\F_q^\times,
  b\in\F_q\}$. The kernel is the group of translations $x\mapsto x+b$, which
  is isomorphic to $(\F_q,+)$. The complement is the group of scalings
  $x\mapsto ax$, which is isomorphic to $\F_q^\times$. Scaling by $a\neq1$
  fixes only $0$, so the action is fixed-point-free. Note that
  $\AGL(1,3)\cong S_3$ and $\AGL(1,4)\cong A_4$.
\end{enumerate}
\end{example}

The next lemma is the special case for abelian groups of Brauer's
permutation lemma. We give the short proof, because it drives the
irreducibility theorem.

\begin{lemma}\label{lem:brauerperm}
Let $\alpha$ be an automorphism of a finite abelian group $A$, and let
$\widehat A=\operatorname{Hom}(A,\C^\times)$. Then $\alpha$ fixes exactly
as many elements of $A$ as characters $\mu\in\widehat A$ with
$\mu\circ\alpha=\mu$. In particular, if $\alpha$ is fixed-point-free then
the trivial character is the only $\alpha$-invariant character.
\end{lemma}

\begin{proof}
Since $A$ is abelian, $\varphi(x)=\alpha(x)x^{-1}$ defines a homomorphism
$\varphi\colon A\to A$. Its kernel is the set of fixed points of $\alpha$. A
character $\mu$ satisfies $\mu\circ\alpha=\mu$ exactly when $\mu$ is
trivial on $\varphi(A)$. The number of such characters is
$|A/\varphi(A)|=|\ker\varphi|$.
\end{proof}

\subsection{Induced representations}\label{sec:induction}

Let $H\le G$ and let $\rho$ be a representation of $H$ with character
$\chi_\rho$. The \emph{induced representation} $\Ind_H^G\rho$ has character
\begin{equation}\label{eq:indchar}
\Ind_H^G\chi_\rho(g)=\frac{1}{|H|}\sum_{\substack{x\in G\\ x^{-1}gx\in H}}
\chi_\rho(x^{-1}gx)
\end{equation}
\cite[\S7.2]{Se77}. A representation induced from a one-dimensional
character is called \emph{monomial}. When $\rho=\triv$ is trivial we write
$\triv_H^G=\Ind_H^G\triv$. This is the \emph{permutation character} of $G$
acting on $G/H$. Its value $\triv_H^G(g)$ is the number of cosets $xH$
fixed by $g$.

We write $\langle\phi,\psi\rangle_G=|G|^{-1}\sum_{g\in G}\phi(g)
\overline{\psi(g)}$. A character $\chi$ is irreducible if and only if
$\langle\chi,\chi\rangle_G=1$. \emph{Frobenius reciprocity} says
$\langle\Ind_H^G\psi,\chi\rangle_G=\langle\psi,\Res_H\chi\rangle_H$.

For $s\in G$ put $H_s=sHs^{-1}\cap H$, and let $\rho^s$ be the
representation $x\mapsto\rho(s^{-1}xs)$ of $H_s$.

\begin{theorem}[Mackey; cited from {\cite[\S7.3--7.4]{Se77}}]
\label{thm:mackey}
\begin{enumerate}[label=(\alph*),itemsep=1pt]
\item Let $s$ run over a set of representatives of the double cosets
  $H\backslash G/H$. Then
  \[
  \langle\Ind_H^G\rho,\Ind_H^G\rho\rangle_G=\sum_{s}
  \langle\rho^s,\Res_{H_s}\rho\rangle_{H_s}.
  \]
\item (Irreducibility criterion.) $\Ind_H^G\rho$ is irreducible if and only
  if $\rho$ is irreducible and, for every $s\in G\setminus H$, the
  representations $\rho^s$ and $\Res_{H_s}\rho$ of $H_s$ have no common
  irreducible constituent.
\end{enumerate}
\end{theorem}

\begin{example}\label{ex:s3ind}
Let $G=S_3$, $H=A_3=\langle c\rangle$ with $c=(123)$, and let $\mu(c)=\omega
=e^{2\pi i/3}$. By \eqref{eq:indchar}, $\Ind_H^G\mu$ takes the value $2$ at
$1$, the value $0$ at transpositions, and the value $\omega+\omega^2=-1$ at
$3$-cycles. Its norm is $(4+0+2\cdot1)/6=1$, so it is the irreducible
$2$-dimensional representation. By contrast, $\triv_{\langle(12)\rangle}^G$
is the permutation character of $S_3$ on three points. It equals $\triv$
plus the $2$-dimensional character.
\end{example}

\subsection{\texorpdfstring{Groups of $2\times2$ matrices}{Groups of 2x2 matrices}}\label{sec:gl2}

Several examples and counterexamples live in the groups $\GL(2,\F_q)$,
$\SL(2,\F_q)$, $\PGL(2,\F_q)=\GL(2,\F_q)/\F_q^\times$ and
$\PSL(2,\F_q)=\SL(2,\F_q)/\{\pm1\}$. Their orders are
\begin{gather*}
|\GL(2,\F_q)|=q(q-1)^2(q+1),\qquad |\SL(2,\F_q)|=|\PGL(2,\F_q)|=q(q^2-1),\\
|\PSL(2,\F_q)|=\frac{q(q^2-1)}{\gcd(2,q-1)}.
\end{gather*}
For even $q$ we have $\PSL(2,\F_q)=\SL(2,\F_q)\cong\PGL(2,\F_q)$, because
$-1=1$ and squaring is bijective on $\F_q^\times$. The group
$\PSL(2,\F_q)$ is simple for $q\ge4$ \cite[Ch.~8]{Ro95}.

A \emph{split torus} of $\GL(2,\F_q)$ is a conjugate of the diagonal
subgroup. A \emph{non-split torus} is a conjugate of $\F_{q^2}^\times$,
acting by multiplication on $\F_{q^2}\cong\F_q^2$. The \emph{unipotent
subgroup} is $U=\{\left(\begin{smallmatrix}1&b\\0&1\end{smallmatrix}\right)
\}$.

\begin{lemma}\label{lem:cyclicvector}
Let $k$ be a field and $g\in M_2(k)$ a non-scalar matrix. Then the
centralizer of $g$ in $M_2(k)$ is the commutative algebra
$k[g]=\{a+bg: a,b\in k\}$.
\end{lemma}

\begin{proof}
Since $g$ is not scalar, some vector $v$ is not an eigenvector. Then $v$ and
$gv$ form a basis. If $x$ commutes with $g$, write $xv=av+bgv$. Then
$x(gv)=g(xv)=(a+bg)gv$. So $x$ and $a+bg$ agree on a basis, and $x=a+bg$.
\end{proof}

\begin{example}\label{ex:smallgl}
The small cases are familiar groups. We have $\GL(2,\F_2)\cong S_3$,
$\PSL(2,\F_3)\cong A_4$ and $\PGL(2,\F_3)\cong S_4$. Also
$\SL(2,\F_4)\cong\PSL(2,\F_5)\cong A_5$. The non-split torus of
$\GL(2,\F_2)$ is $\F_4^\times$, cyclic of order~$3$. Under $\GL(2,\F_2)\cong
S_3$ it corresponds to $A_3$.
\end{example}

\subsection{Graphs and harmonic morphisms}\label{sec:harmbg}

Harmonic morphisms are the graph-theoretic analogue of holomorphic maps
between Riemann surfaces \cite{BN09,Ur00}. They are treated at length in
\cite{MJ26}. This subsection gives the definitions we need.

All graphs here are finite and simple, with no loops or multiple edges. The
\emph{genus} of a connected graph $\Gamma$ is $g(\Gamma)=|E(\Gamma)|-
|V(\Gamma)|+1$, its first Betti number. For example
$g(K_n)=\binom n2-n+1=(n-1)(n-2)/2$.

\begin{definition}[Baker--Norine \cite{BN09}]\label{def:harmonic}
Let $\Gamma$ and $\Gamma'$ be graphs. A \emph{morphism}
$\varphi\colon\Gamma\to\Gamma'$ sends vertices to vertices and edges to
edges or vertices, as follows. For each edge $e=\{x,y\}$ of $\Gamma$,
either $\varphi(x)\neq\varphi(y)$ are adjacent and
$\varphi(e)=\{\varphi(x),\varphi(y)\}$, or $\varphi(x)=\varphi(y)$ and
$\varphi(e)=\varphi(x)$. In the second case $e$ is \emph{vertical}. Let
$v_\varphi(x)$ be the number of vertical edges at $x$. For a vertex $x$ and
an edge $e'$ of $\Gamma'$ at $\varphi(x)$, let $m_\varphi(x,e')$ be the
number of edges $e$ at $x$ with $\varphi(e)=e'$. The morphism is
\emph{harmonic at $x$} if $m_\varphi(x,e')$ takes the same value
$m_\varphi(x)$ for every edge $e'$ at $\varphi(x)$. It is \emph{harmonic}
if it is harmonic at every vertex. The number $m_\varphi(x)$ is the
\emph{horizontal multiplicity} at $x$.
\end{definition}

If $\varphi(x)$ has no edges, the condition at $x$ is vacuous and we put
$m_\varphi(x)=0$. Our graphs are often disconnected. Harmonicity is a
condition at each vertex, so we may check it one component at a time.

\begin{theorem}[Baker--Norine; cited from {\cite[\S2]{BN09}}]
\label{thm:bn}
Let $\varphi\colon\Gamma\to\Gamma'$ be a harmonic morphism of connected
graphs, where $\Gamma'$ has at least one edge and $m_\varphi(x)\ge1$ for
all $x$.
\begin{enumerate}[label=(\alph*),itemsep=1pt]
\item The number of edges over an edge $e'$ of $\Gamma'$ does not depend on
  $e'$. It is the \emph{degree} $\deg\varphi$. For every vertex $y$ of
  $\Gamma'$, $\deg\varphi=\sum_{\varphi(x)=y}m_\varphi(x)$.
\item (Riemann--Hurwitz.) $2g(\Gamma)-2=\deg\varphi\,\bigl(2g(\Gamma')-2
  \bigr)+\sum_{x}\bigl(2(m_\varphi(x)-1)+v_\varphi(x)\bigr)$.
\end{enumerate}
\end{theorem}

Two families of maps occur below. In each, the vertices of a graph are
partitioned into blocks and each block is collapsed to one vertex.

\begin{lemma}\label{lem:collapse}
\begin{enumerate}[label=(\alph*),itemsep=1pt]
\item Partition the vertices of $K_n$ into $k\ge2$ blocks of sizes
  $a_1,\dots,a_k$. Let $\varphi\colon K_n\to K_k$ collapse each block to a
  vertex. Then $\varphi$ is harmonic if and only if $k=2$ or all $a_i$ are
  equal. If all $a_i=a$, then $m_\varphi(x)=a$ and $v_\varphi(x)=a-1$ for
  every $x$, and $\deg\varphi=a^2$.
\item Let $\Gamma$ be the complete multipartite graph with parts
  $P_1,\dots,P_M$, where $M\ge2$. Let $\psi\colon\Gamma\to K_M$ collapse
  each part to a vertex. Then $\psi$ is harmonic if and only if $M=2$ or all
  $|P_j|$ are equal. If all $|P_j|=s$, then $m_\psi(x)=s$ and
  $v_\psi(x)=0$ for every $x$, and $\deg\psi=s^2$.
\end{enumerate}
In both cases, when the blocks have equal size, the two sides of the
Riemann--Hurwitz formula agree.
\end{lemma}

\begin{proof}
(a) Let $x$ lie in block $j$. The edges at $x$ over the edge $\{j,l\}$ of
$K_k$ are the edges from $x$ to block $l$, and there are $a_l$ of them. So
$\varphi$ is harmonic at $x$ if and only if $a_l$ is the same for all $l\neq
j$. If $k=2$ this is automatic. If $k\ge3$, the conditions for $j=1$ and
$j=2$ force $a_2=\dots=a_k$ and $a_1=a_3$, so all blocks have equal size.
The vertical edges at $x$ join $x$ to the rest of its block. Each edge of
$K_k$ has $a\cdot a$ edges over it.

(b) The argument is the same. A part is an independent set, so no edge is
vertical. The edges at $x\in P_i$ over $\{i,j\}$ are the $|P_j|$ edges from
$x$ to $P_j$.

For Riemann--Hurwitz in (a), put $n=ak$. Then $2g(K_n)-2=n(n-3)$, and
\[
a^2k(k-3)+n\bigl(2(a-1)+(a-1)\bigr)=n^2-3an+3an-3n=n(n-3).
\]
In (b), put $n=Ms$. Then $\Gamma$ has $s^2M(M-1)/2$ edges, so
$2g(\Gamma)-2=s^2M(M-1)-2Ms$, and
\[
s^2(M^2-3M)+Ms\cdot2(s-1)=s^2M^2-s^2M-2Ms.\qedhere
\]
\end{proof}

\begin{example}\label{ex:harmbg}
Collapse $K_4$ onto $K_2$ by $1,2\mapsto u$ and $3,4\mapsto w$. The edges
$\{1,2\}$ and $\{3,4\}$ are vertical. The other four edges lie over
$\{u,w\}$, so the map is harmonic with $m=2$ at every vertex and degree~$4$.
Riemann--Hurwitz reads $4=4\cdot(-2)+4\cdot3$, since $g(K_4)=3$ and
$g(K_2)=0$.

Now collapse $K_4$ onto $K_3$ by $1,2\mapsto u$, $3\mapsto w$ and $4\mapsto
z$. At vertex $3$ there are two edges over $\{w,u\}$ but only one over
$\{w,z\}$. So this map is not harmonic. The blocks have unequal sizes, as
Lemma~\ref{lem:collapse}(a) predicts.
\end{example}

%======================================================================
\section{Notation}\label{sec:notation}

\begin{center}
\renewcommand{\arraystretch}{1.15}
\begin{tabular}{>{$}l<{$}p{11.5cm}}
\toprule
\text{Symbol} & Meaning\\
\midrule
Z=Z(G) & center of $G$\\
\Cent(g,G) & centralizer of $g$ in $G$\\
N_G(H) & normalizer of the subgroup $H$\\
W(C) & the quotient $N_G(C)/C$ (Section~\ref{sec:induced})\\
\Conj(g,G) & conjugacy class $\{x^{-1}gx : x\in G\}$\\
G_* & set of conjugacy classes of $G$\\
\widehat C & group of characters $C\to\C^\times$ of an abelian group $C$\\
\Ind_H^G,\ \Res_H & induction and restriction\\
\triv,\ \triv_H^G & trivial character and permutation character
$\Ind_H^G\triv$\\
\langle\ ,\ \rangle_G & inner product of class functions on $G$\\
\rho^s & conjugate representation $x\mapsto\rho(s^{-1}xs)$ on
$sHs^{-1}\cap H$\\
R(G) & character ring: $\Z$-linear combinations of irreducible characters\\
\Gamma(G) & commuting graph on $G\setminus Z$ (Section~\ref{sec:graph})\\
\Gamma^c(G) & non-commuting graph, the complement of $\Gamma(G)$\\
\Delta(G) & commuting conjugacy class graph (Definition~\ref{def:ccc})\\
\varphi_G & class map $\Gamma(G)\to\Delta(G)$, $x\mapsto$ class of $x$\\
m_\varphi(x),\ v_\varphi(x) & horizontal multiplicity and number of vertical
edges at $x$ (Definition~\ref{def:harmonic})\\
g(\Gamma) & genus $|E|-|V|+1$ of a connected graph\\
\supp(H),\ \tau(H) & support of $H$ in $G_*$ and its indicator vector in
$\F_2^{G_*}$ (Section~\ref{sec:perm})\\
\F_q & finite field with $q$ elements\\
T,\ T',\ U & split torus, non-split torus, unipotent subgroup\\
\phi & golden ratio $(1+\sqrt5)/2$\\
\bottomrule
\end{tabular}
\end{center}

%======================================================================
\section{Commutative transitive and commutation transitive groups}
\label{sec:ct}

This section defines the two classes of groups, gives examples and
non-examples, and recalls the classification. It ends with three false
claims from the 2000 draft, each refuted.

\subsection{Definitions and first examples}\label{sec:defs}

\begin{definition}\label{def:ct}
A group $G$ is \emph{commutative transitive}, or CT, if for all non-identity
$x,y,w\in G$, the equations $xy=yx$ and $yw=wy$ imply $xw=wx$. It is
\emph{commutation transitive} if the same holds for all non-central
$x,y,w$.
\end{definition}

CT groups are also called CA-groups. Commutation transitive groups are also
called AC-groups (see Section~\ref{sec:lit}). We avoid the abbreviations
except ``CT'', to keep the two notions apart.

\begin{proposition}\label{prop:basic}
Let $G$ be a group.
\begin{enumerate}[label=(\alph*),itemsep=1pt]
\item $G$ is commutation transitive if and only if $\Cent(g,G)$ is abelian
  for every $g\notin Z$.
\item $G$ is CT if and only if $\Cent(g,G)$ is abelian for every $g\neq1$.
\item If $G$ is CT and non-abelian, then $Z=1$.
\item Every CT group is commutation transitive. A commutation transitive
  group with $Z=1$ is CT.
\end{enumerate}
\end{proposition}

\begin{proof}
(a) Suppose $G$ is commutation transitive, and let $g\notin Z$ and
$x,y\in\Cent(g,G)$. If $x$ or $y$ is central, they commute. Otherwise $x$
commutes with $g$ and $g$ commutes with $y$, so $x$ commutes with $y$.
Conversely, suppose all these centralizers are abelian, and let $x,y,w$ be
non-central with $xy=yx$ and $yw=wy$. Then $x,w\in\Cent(y,G)$, which is
abelian, so $xw=wx$.

(b) The proof is the same, with ``non-central'' replaced by
``non-identity''.

(c) Let $1\neq z\in Z$. Then $\Cent(z,G)=G$, which is abelian by (b).

(d) The first claim follows from (a) and (b). The second holds because
``non-central'' and ``non-identity'' mean the same thing when $Z=1$.
\end{proof}

\begin{example}\label{ex:ctexamples}
These groups are CT.
\begin{enumerate}[itemsep=1pt]
\item Every abelian group.
\item $S_3$. The centralizers are $A_3$ and the three subgroups of order
  two.
\item The Frobenius groups of Example~\ref{ex:frob}. This follows from
  Proposition~\ref{prop:frobct} below.
\item $A_5\cong\SL(2,\F_4)$ (Proposition~\ref{prop:psl2even}).
\item Free groups. A centralizer of a non-trivial element of a free group is
  infinite cyclic \cite{LR86,FGRS15}. (Cited.)
\end{enumerate}
\end{example}

\begin{example}\label{ex:nonct}
The quaternion group $Q_8$ and the dihedral group $D_8$ of order~$8$ are
commutation transitive but not CT. Each has center of order~$2$. In $Q_8$
the centralizer of $i$ is $\langle i\rangle$. In $D_8=\langle r,s\rangle$
the centralizer of $r$ is $\langle r\rangle$, and that of $s$ is
$\{1,s,r^2,sr^2\}$. All of these are abelian. By
Proposition~\ref{prop:basic}(c), neither group is CT.

The group $S_4$ is neither. Its center is trivial, and
$\Cent((12)(34),S_4)$ is dihedral of order~$8$.
\end{example}

\begin{proposition}\label{prop:gl2}
For every field $k$, the groups $\GL(2,k)$ and $\SL(2,k)$ are commutation
transitive.
\end{proposition}

\begin{proof}
The centers consist of scalar matrices. If $g$ is not central it is not
scalar. By Lemma~\ref{lem:cyclicvector}, its centralizer is $k[g]^\times$
or $k[g]^\times\cap\SL(2,k)$. Both are abelian. Now apply
Proposition~\ref{prop:basic}(a).
\end{proof}

\subsection{The commuting graph}\label{sec:graph}

The \emph{commuting graph} $\Gamma(G)$ has vertex set $G\setminus Z$. Two
distinct vertices are joined by an edge when they commute.

\begin{proposition}\label{prop:graph}
Let $G$ be commutation transitive. Each connected component of $\Gamma(G)$
is a complete graph. The map $K\mapsto K\cup Z$ is a bijection from the set
of components to the set $\{\Cent(g,G): g\notin Z\}$. Its inverse is
$C\mapsto C\setminus Z$.
\end{proposition}

\begin{proof}
Transitivity of commuting on $G\setminus Z$ makes ``equal or adjacent'' an
equivalence relation. Its classes are the components, and any two vertices
in a class are adjacent. The class of $g$ is $\Cent(g,G)\setminus Z$, which
proves the bijection.
\end{proof}

\begin{example}\label{ex:a5graph}
For $A_5$ the graph has $59$ vertices. There are $5$ triangles, one for each
Klein four-group: $5\cdot3=15$ involutions. There are $10$ edges, one for
each subgroup of order~$3$: $20$ elements of order~$3$. And there are $6$
copies of $K_4$, one for each subgroup of order~$5$: $24$ elements of
order~$5$. Section~\ref{sec:harmonic} studies maps out of this graph.
\end{example}

\subsection{Centralizers are malnormal}\label{sec:malnormal}

\begin{lemma}\label{lem:cent}
Let $G$ be commutation transitive, $g_0\notin Z$ and $C=\Cent(g_0,G)$.
\begin{enumerate}[label=(\alph*),itemsep=1pt]
\item If $g_1\in C\setminus Z$, then $\Cent(g_1,G)=C$.
\item The centralizer of the subgroup $C$ is $C$ itself.
\item For every $x\in G$, either $xCx^{-1}=C$, or $C\cap xCx^{-1}=Z$.
\end{enumerate}
\end{lemma}

\begin{proof}
(a) Since $C$ is abelian, $C\subseteq\Cent(g_1,G)$. Conversely, if $y$
commutes with $g_1$ and $y\notin Z$, then $y$ commutes with $g_0$ by
transitivity. Central $y$ lie in $C$ anyway.

(b) Since $C$ is abelian, it centralizes itself. Anything centralizing $C$
commutes with $g_0$, so it lies in $C$.

(c) The set $C\cap xCx^{-1}$ contains $Z$. Suppose it contains some
$c\notin Z$. By Lemma~\ref{lem:conjcent}, $xCx^{-1}=\Cent(xg_0x^{-1},G)$,
and $xg_0x^{-1}\notin Z$. By (a), applied to both centralizers,
$C=\Cent(c,G)=xCx^{-1}$.
\end{proof}

The 2000 draft states (c) as ``two centralizers are disjoint or equal''.
That is false as written, since every subgroup contains~$1$ and $Z$.

\begin{proposition}\label{prop:selfnorm}
Let $G$ be commutation transitive, $g_0\notin Z$ and $C=\Cent(g_0,G)$.
Suppose $N_G(C)=C$. Then $C/Z$ is a Frobenius complement in $G/Z$. In
particular, $G/Z$ has a proper non-trivial normal subgroup, its Frobenius
kernel.
\end{proposition}

\begin{proof}
We have $Z\subsetneq C\subsetneq G$, because $g_0\in C\setminus Z$ and $C$
is abelian while $G$ is not. Let $\bar x\in G/Z$ lie outside $C/Z$. Then
$x\notin N_G(C)$, so $C\cap xCx^{-1}=Z$ by Lemma~\ref{lem:cent}(c). Both
subgroups contain $Z$, so $(C/Z)\cap\bar x(C/Z)\bar x^{-1}=
(C\cap xCx^{-1})/Z=1$. Now apply Theorem~\ref{thm:frobenius}.
\end{proof}

\begin{example}\label{ex:selfnorm}
In $S_3$ the centralizer $\langle(12)\rangle$ is self-normalizing. Its
Frobenius kernel is $A_3$. In $A_5$, Sylow counting gives the normalizers.
There are $6$ Sylow $5$-subgroups, $10$ Sylow $3$-subgroups and $5$ Sylow
$2$-subgroups. So the centralizers $C_5$, $C_3$ and $V_4$ have normalizers
of orders $10$, $6$ and $12$. None is self-normalizing, as
Proposition~\ref{prop:selfnorm} requires for a simple group.
\end{example}

\subsection{The classification of finite CT groups}\label{sec:classif}

\begin{theorem}[Weisner, Suzuki, Brauer--Suzuki--Wall, Wu. Cited]
\label{thm:classif}
A finite group $G$ is CT if and only if one of the following holds.
\begin{enumerate}[label=(\roman*),itemsep=1pt]
\item $G$ is abelian.
\item $G=N\rtimes H$ is a Frobenius group with abelian kernel $N$ and cyclic
  complement $H\neq1$.
\item $G\cong\PSL(2,\F_{2^f})$ for some $f\ge2$.
\end{enumerate}
\end{theorem}

The sources are as follows. Weisner showed that a finite CT group is
solvable or simple \cite{We25}. Wu gave a complete proof, and described the
solvable case as the semidirect product of the Fitting subgroup by a
fixed-point-free cyclic group \cite{Wu98}. Suzuki and Brauer--Suzuki--Wall
identified the simple ones \cite{Su57,BSW58}. A cyclic group of prime order
is abelian, so it falls under~(i).

\begin{remark}\label{rem:weisnerk}
Weisner states the simple case with $f>2$. But $\PSL(2,\F_4)$ has order
$60$, so it is isomorphic to $A_5$, which is CT. This was already noted in
the 2000 draft.
\end{remark}

Both directions of ``if'' can be checked directly.

\begin{proposition}\label{prop:frobct}
Let $N$ be a finite abelian group and $H\neq1$ a cyclic group acting
fixed-point-freely on $N$. Then $G=N\rtimes H$ is CT. Its centralizers of
non-identity elements are $N$ and the conjugates of $H$. Moreover $H$ is a
Frobenius complement with kernel $N$.
\end{proposition}

\begin{proof}
Let $1\neq n\in N$. If $n'h$ commutes with $n$ (with $n'\in N$, $h\in H$),
then $hnh^{-1}=n$, because $N$ is abelian. So $h=1$, and $\Cent(n,G)=N$.

Now let $g=n'h$ with $h\neq1$. Conjugation by $g$ on $N$ equals conjugation
by $h$, which is fixed-point-free. So $\Cent(g,G)\cap N=1$. Hence
$\Cent(g,G)$ maps injectively into $G/N\cong H$, and it is cyclic.

Next, $H$ is malnormal. Let $x=nh'\notin H$, so that $n\neq1$ and
$xHx^{-1}=nHn^{-1}$. If $h\in H\cap nHn^{-1}$, write $h=nh_2n^{-1}$.
Reducing modulo $N$ gives $h=h_2$, so $h$ commutes with $n$, and therefore
$h=1$. Finally, $g$ lies in a conjugate of $H$. The conjugates of $H$ cover
$|N|(|H|-1)+1$ elements, and the remaining elements form $N$ by
Theorem~\ref{thm:frobenius}. So $H\subseteq\Cent(g,G)$ up to conjugation,
and the orders force equality.
\end{proof}

\begin{proposition}\label{prop:psl2even}
Let $q=2^f\ge4$ and $G=\SL(2,\F_q)=\PSL(2,\F_q)$. Every non-identity
element of $G$ is of exactly one of three kinds.
\begin{enumerate}[label=(\roman*),itemsep=1pt]
\item Unipotent. Its centralizer is a conjugate of $U\cong(\F_q,+)$, of
  order $q$.
\item Split semisimple. Its centralizer is a conjugate of
  $T=\{\diag(a,a^{-1})\}\cong\F_q^\times$, of order $q-1$.
\item Non-split semisimple. Its centralizer is a conjugate of $T'$, the
  norm-one subgroup of $\F_{q^2}^\times$, cyclic of order $q+1$.
\end{enumerate}
In particular $G$ is CT. There are $q+1$ conjugacy classes. One is $\{1\}$,
one is unipotent, $(q-2)/2$ are split and $q/2$ are non-split. The
normalizers $N_G(U)$, $N_G(T)$, $N_G(T')$ have orders $q(q-1)$, $2(q-1)$ and
$2(q+1)$.
\end{proposition}

\begin{proof}
Let $g\neq1$ with trace $t$. Its characteristic polynomial is $x^2+tx+1$.

If $t=0$, this polynomial is $(x+1)^2$ and $g$ is a non-trivial unipotent
element. It is conjugate to $u=\left(\begin{smallmatrix}1&1\\0&1
\end{smallmatrix}\right)$. A direct computation gives
$\Cent(u,G)=\{\left(\begin{smallmatrix}a&b\\0&a\end{smallmatrix}\right):
a^2=1\}=U$.

If $t\neq0$ and the polynomial has roots $a,a^{-1}\in\F_q$, then $a\neq
a^{-1}$, because $a^2=1$ forces $a=1$ in characteristic $2$. So $g$ is
conjugate to $\diag(a,a^{-1})$. By Lemma~\ref{lem:cyclicvector}, its
centralizer is $T$.

If the polynomial is irreducible, then $\F_q[g]\cong\F_{q^2}$, and the
centralizer is the norm-one subgroup $T'$.

All three centralizers are abelian, so $G$ is CT by
Proposition~\ref{prop:basic}(b).

For the class count, the transvections $\left(\begin{smallmatrix}1&b\\0&1
\end{smallmatrix}\right)$ with $b\ne0$ are all conjugate, since conjugating
by $\diag(c,c^{-1})$ multiplies $b$ by $c^2$ and squaring is onto. A split
element is determined up to conjugacy by the pair $\{a,a^{-1}\}$, which
gives $(q-2)/2$ classes. A non-split element is determined by a pair
$\{z,z^q\}=\{z,z^{-1}\}$ with $1\ne z\in T'$, which gives $q/2$ classes.

For the normalizers, $U$ is a Sylow $2$-subgroup. Its conjugates correspond
to the $q+1$ points of the projective line. Split tori correspond to the
$q(q+1)/2$ unordered pairs of distinct points, and the non-split tori number
$q(q-1)/2$. Dividing $|G|=q(q^2-1)$ by these numbers gives the stated
orders. As a check, these subgroups partition the non-identity elements:
\[
(q+1)(q-1)+\tfrac{q(q+1)}{2}(q-2)+\tfrac{q(q-1)}{2}\,q=q^3-q-1=|G|-1.
\qedhere
\]
\end{proof}

\begin{example}\label{ex:a5classes}
For $q=4$ we get $A_5$. There are five classes, of sizes $1,15,20,12,12$.
The centralizers are $V_4$, $C_3$ and $C_5$, as in
Example~\ref{ex:a5graph}. The two classes of $5$-cycles are the classes of
$a$ and $a^2$, for a generator $a$ of $C_5$. We call them $5A$ and $5B$.
\end{example}

\subsection{Three false claims of the 2000 draft}\label{sec:false}

\begin{counterexample}[$\PGL(2,\F_q)$ is not CT]\label{cex:pgl}
The 2000 draft claimed that $\PGL(2,\F_q)$ is CT for odd $q$. This is false
for every odd $q$. Let $t$ be the image of $\diag(-1,1)$ and $w$ the image
of $\left(\begin{smallmatrix}0&1\\1&0\end{smallmatrix}\right)$. Conjugation
by $w$ sends $\diag(-1,1)$ to $\diag(1,-1)=-\diag(-1,1)$, so $w$ commutes
with $t$ in $\PGL(2,\F_q)$. The image $\bar T$ of the diagonal torus also
commutes with $t$. But $w$ inverts $\bar T$, which is cyclic of order
$q-1$. So $\Cent(t)\supseteq\langle\bar T,w\rangle$, which is dihedral of
order $2(q-1)$ and non-abelian when $q\ge5$. For $q=3$ the group is $S_4$,
which is not CT by Example~\ref{ex:nonct}. The center of $\PGL(2,\F_q)$ is
trivial, so the group is not commutation transitive either.

Ironically, the draft's own footnote supplies this pair $t,w$.
\end{counterexample}

The same idea settles the odd case of $\PSL(2,\F_q)$.

\begin{proposition}\label{prop:psl2odd}
For odd $q$, $\PSL(2,\F_q)$ is CT if and only if $q\in\{3,5\}$.
\end{proposition}

\begin{proof}
For $q=3,5$ the group is $A_4$ or $A_5$, which are CT
(Example~\ref{ex:ctexamples}). Let $q\ge7$, and write bars for images in
$\PSL(2,\F_q)$.

Suppose first $q\equiv1\pmod4$. Choose $i\in\F_q$ with $i^2=-1$, let
$\tilde t=\diag(i,-i)$ and $w=\left(\begin{smallmatrix}0&1\\-1&0
\end{smallmatrix}\right)$. Then $w\tilde tw^{-1}=-\tilde t$, so $\bar w$
commutes with $\bar t$. The image $\bar T$ of $\{\diag(a,a^{-1})\}$ is
abelian, contains $\bar t$, and has order $(q-1)/2$. Conjugation by $w$
sends $\diag(a,a^{-1})$ to its inverse. So $\bar w$ commutes with
$\diag(a,a^{-1})$ only if $a^{-1}=\pm a$, that is $a^4=1$. At most two
elements of $\bar T$ satisfy this. Since $|\bar T|\ge4$ when $q\ge9$,
$\Cent(\bar t)$ is non-abelian.

Suppose next $q\equiv3\pmod4$. Let $\tilde t=\left(\begin{smallmatrix}0&1\\
-1&0\end{smallmatrix}\right)$. Since $\tilde t^2=-1$ and $-1$ is not a
square, $\F_q[\tilde t]\cong\F_{q^2}$. Its norm-one group $S$ is cyclic of
order $q+1$, and $\bar S$ is abelian of order $(q+1)/2\ge4$. Every element
of $\F_q$ is a sum of two squares, so we can choose $a,b$ with
$-a^2-b^2=1$. Then $x=\left(\begin{smallmatrix}a&b\\b&-a\end{smallmatrix}
\right)$ lies in $\SL(2,\F_q)$ and satisfies $x\tilde t=-\tilde tx$. Hence
$x$ acts on $\F_q[\tilde t]$ as the non-trivial field automorphism. That
automorphism inverts $S$. As before, at most two elements of $\bar S$
commute with $\bar x$. So $\Cent(\bar t)\supseteq\langle\bar S,\bar x
\rangle$ is non-abelian.
\end{proof}

The accompanying script confirms that $\PSL(2,\F_7)$ is not CT, while
$\PSL(2,\F_3)$ and $\PSL(2,\F_5)$ are.

\begin{counterexample}[Commutation in $G/Z$ is not commutation in $G$]
\label{cex:quotient}
The 2000 draft transported the relation of commuting from $G\setminus Z$ to
$(G\setminus Z)/Z$, by declaring $\bar x\sim\bar y$ when $x$ and $y$
commute. That relation is transitive. But it is not the same as commuting in
$G/Z$. For odd $q$, $\GL(2,\F_q)$ is commutation transitive
(Proposition~\ref{prop:gl2}). Its quotient by the center is
$\PGL(2,\F_q)$, which is not (Counterexample~\ref{cex:pgl}). So the class of
commutation transitive groups is not closed under passing to $G/Z(G)$. By
contrast, Wu shows that the class of locally finite CT groups is closed
under quotients \cite{Wu98}. (Cited.)
\end{counterexample}

%======================================================================
\section{Inducing characters from centralizers}\label{sec:induced}

This section computes $\Ind_C^G\mu$ when $C$ is a centralizer, and decides
exactly when it is irreducible.

Throughout, $G$ is a finite commutation transitive group, $g_0\notin Z$,
$C=\Cent(g_0,G)$, and $\mu\in\widehat C$. Since $C$ is abelian, $\mu$ is
one-dimensional. The group $N_G(C)$ acts on $C$ by conjugation. Since $C$
acts trivially, this is an action of $W(C)=N_G(C)/C$.

\subsection{The character formula}

\begin{proposition}\label{prop:indformula}
Let $\chi=\Ind_C^G\mu$ and $g\in G$.
\begin{enumerate}[label=(\alph*),itemsep=1pt]
\item If $g\in Z$, then $\chi(g)=[G:C]\,\mu(g)$.
\item If $g\notin Z$ and $\Conj(g,G)\cap C=\emptyset$, then $\chi(g)=0$.
\item If $g\notin Z$ and $a\in\Conj(g,G)\cap C$, then $\Conj(g,G)\cap C$ is
  the $W(C)$-orbit of $a$, and
  \begin{equation}\label{eq:indformula}
  \chi(g)=\sum_{n\in W(C)}\mu(n^{-1}an)
  =\frac{1}{|C|}\sum_{n\in N_G(C)}\mu(n^{-1}an).
  \end{equation}
\end{enumerate}
\end{proposition}

\begin{proof}
Use \eqref{eq:indchar}. In (a), $x^{-1}gx=g\in C$ for all $x$, so all
$|G|$ terms equal $\mu(g)$. Part (b) is immediate. For (c), write
$g=sas^{-1}$ and put $u=s^{-1}x$. Then $x^{-1}gx=u^{-1}au$. Since $a\in
C\setminus Z$, Lemma~\ref{lem:cent}(a) gives $\Cent(a,G)=C$. If
$u^{-1}au\in C$, then $u^{-1}au\in C\setminus Z$, and
\[
C=\Cent(u^{-1}au,G)=u^{-1}\Cent(a,G)u=u^{-1}Cu
\]
by Lemma~\ref{lem:conjcent}. So $u\in N_G(C)$. Conversely, every $u\in
N_G(C)$ satisfies $u^{-1}au\in C$. Thus the $x$ in \eqref{eq:indchar} form
the coset $sN_G(C)$, and the elements $u^{-1}au$ form the $W(C)$-orbit of
$a$. The term $\mu(u^{-1}au)$ depends only on $uC$, because $C$ is abelian.
\end{proof}

\begin{counterexample}[The 2000 formula]\label{cex:formula}
The 2000 draft stated \eqref{eq:indformula} with the sum over $N_G(C)$ and
without the factor $|C|^{-1}$. It also gave $|N_G(C)|\mu(g)$ for central
$g\neq1$, instead of (a). Take $G=S_3$, $C=A_3$ and $\mu$ non-trivial, as in
Example~\ref{ex:s3ind}. At a $3$-cycle the true value is $-1$, while the
draft's formula gives $\omega+\omega^2$ summed over $N_G(C)=S_3$, which is
$-3$.
\end{counterexample}

\begin{example}\label{ex:a5ind}
Let $G=A_5$ and $C=\langle a\rangle\cong C_5$, with $\mu_k(a)=\zeta^k$ and
$\zeta=e^{2\pi i/5}$. Here $|W(C)|=10/5=2$, and $W(C)$ acts by $a\mapsto
a^{-1}$. By \eqref{eq:indformula}, $\Ind_C^G\mu_1$ is $12$ at $1$, zero on
the classes of orders $2$ and $3$, $\zeta+\zeta^{-1}=\phi-1$ on $5A$ and
$\zeta^2+\zeta^{-2}=-\phi$ on $5B$. Here $\phi=(1+\sqrt5)/2$.
\end{example}

\subsection{Irreducibility}

\begin{theorem}\label{thm:irred}
Let $G$ be a finite non-abelian CT group, $1\neq g_0\in G$,
$C=\Cent(g_0,G)$ and $\mu\in\widehat C$. Then $\Ind_C^G\mu$ is irreducible
if and only if $C$ is normal in $G$ and $\mu\neq\triv$.
\end{theorem}

\begin{proof}
Here $Z=1$ by Proposition~\ref{prop:basic}(c), and $C\neq G$ because $C$ is
abelian. We use Mackey's criterion, Theorem~\ref{thm:mackey}(b).

Suppose $C$ is not normal, and choose $s\notin N_G(C)$. By
Lemma~\ref{lem:cent}(c), $sCs^{-1}\cap C=1$. The restrictions of $\mu^s$ and
$\mu$ to the trivial group are both the trivial representation. So the
criterion fails, and $\Ind_C^G\mu$ is reducible.

Suppose $C$ is normal, and let $s\in G\setminus C$. Then $sCs^{-1}\cap C=C$
and $\mu^s=\mu\circ\alpha$, where $\alpha(c)=s^{-1}cs$. This automorphism
of $C$ is fixed-point-free. Indeed, if $1\neq c\in C$ commutes with $s$,
then $s\in\Cent(c,G)=C$ by Lemma~\ref{lem:cent}(a), which is a
contradiction. By Lemma~\ref{lem:brauerperm}, $\mu^s=\mu$ holds only for
$\mu=\triv$. So for $\mu\neq\triv$ the criterion holds, and $\Ind_C^G\mu$
is irreducible. For $\mu=\triv$, Frobenius reciprocity shows that
$\Ind_C^G\triv$ contains $\triv$. Its degree $[G:C]$ is larger than $1$, so
it is reducible.
\end{proof}

\begin{corollary}\label{cor:frobchars}
\begin{enumerate}[label=(\alph*),itemsep=1pt]
\item If $G\cong\PSL(2,\F_{2^f})$ with $f\ge2$, no character induced from a
  centralizer is irreducible.
\item Let $G=N\rtimes H$ be as in Theorem~\ref{thm:classif}(ii). The
  irreducible characters of $G$ are the $|H|$ linear characters inflated
  from $G/N\cong H$, together with the characters $\Ind_N^G\mu$. Here $\mu$
  runs over representatives of the $H$-orbits on $\widehat N\setminus
  \{\triv\}$. There are $(|N|-1)/|H|$ such orbits, and each of these
  characters has degree $|H|$.
\end{enumerate}
\end{corollary}

\begin{proof}
(a) A simple group has no proper non-trivial normal subgroup, and
centralizers are proper and non-trivial.

(b) The characters $\Ind_N^G\mu$ with $\mu\ne\triv$ are irreducible by
Theorem~\ref{thm:irred}, since $N=\Cent(n,G)$ is normal
(Proposition~\ref{prop:frobct}). By Lemma~\ref{lem:brauerperm}, $H$ acts
freely on $\widehat N\setminus\{\triv\}$. So there are $(|N|-1)/|H|$
orbits. Since $N$ is normal, $\Res_N\Ind_N^G\mu=\sum_{h\in H}\mu^h$. By
Frobenius reciprocity, $\langle\Ind_N^G\mu,\Ind_N^G\mu'\rangle$ is the
number of $h$ with $\mu^h=\mu'$. So characters from different orbits are
distinct. The degrees satisfy
\[
|H|\cdot1^2+\frac{|N|-1}{|H|}\,|H|^2=|N||H|=|G|.
\]
Hence the list is complete.
\end{proof}

\begin{example}\label{ex:frobchars}
For $A_4=V_4\rtimes C_3$, Corollary~\ref{cor:frobchars} gives three linear
characters and one character of degree~$3$, namely $\Ind_{V_4}^{A_4}\lambda$
for any $\lambda\ne\triv$. For $\AGL(1,q)$ it gives $q-1$ linear characters
and one irreducible character of degree $q-1$. The script verifies
Theorem~\ref{thm:irred} exhaustively for $S_3$, $A_4$ and $A_5$, for every
centralizer and every character.
\end{example}

\subsection{\texorpdfstring{Tori in $\PGL(2,\F_q)$}{Tori in PGL(2,Fq)}}

The 2000 draft also stated a reducibility result for $\PGL(2,\F_q)$. Since
that group is not CT for odd $q$, the results above do not apply to it. We
compute the answer directly with Mackey's formula.

\begin{proposition}\label{prop:pglnorm}
Let $G=\PGL(2,\F_q)$. Let $T\le G$ be the image of a maximal torus of
$\GL(2,\F_q)$, either split (with $q\ge3$) or non-split. Let $\pi\in
\widehat T$, let $w$ be a generator of $N_G(T)/T$, and put $\delta=1$ if
$\pi^w=\pi$ and $\delta=0$ otherwise. Then
\[
\langle\Ind_T^G\pi,\Ind_T^G\pi\rangle_G=
\begin{cases}
q+3+\delta, & T \text{ split},\\
q-1+\delta, & T \text{ non-split}.
\end{cases}
\]
\end{proposition}

\begin{proof}
Let $\tilde T\le\GL(2,\F_q)$ be the torus and $A=\F_q[\tilde T]$ its
$\F_q$-span. This is $\F_q\times\F_q$ in the split case and $\F_{q^2}$ in
the non-split case. The hypothesis $q\ge3$ in the split case ensures that
$\tilde T$ spans~$A$.

\emph{Step 1: $T$ is cyclic, and $[N_G(T):T]=2$.} We have $T\cong
\F_q^\times$ (via $\diag(a,b)\mapsto a/b$) or $T\cong\F_{q^2}^\times/
\F_q^\times$. Both are cyclic, of orders $q-1$ and $q+1$. The preimage of
$T$ is $\tilde T$, because $\tilde T$ contains the scalars. So the preimage
of $N_G(T)$ is $N(\tilde T)$, which normalizes $A$. By
Lemma~\ref{lem:cyclicvector} the kernel of the action of $N(\tilde T)$ on
$A$ is $\tilde T$. The image is the automorphism group of $A$ over $\F_q$,
which has order~$2$. The non-trivial automorphism is realized by the swap
matrix in the split case, and by the Frobenius map $z\mapsto z^q$ of
$\F_{q^2}$ in the non-split case. The non-trivial
automorphism acts on $T$ by inversion. So $N_G(T)$ is dihedral of order
$2|T|$, and its elements outside $T$ are involutions.

\emph{Step 2: if $s\notin N_G(T)$ then $T\cap sTs^{-1}=1$.} Suppose $1\neq
t\in T\cap sTs^{-1}$, and lift $t$ to a non-scalar $\tilde t\in\tilde T$.
If $x\in\GL(2,\F_q)$ maps to $\Cent(t,G)$, then $x\tilde tx^{-1}=\lambda
\tilde t$ for some $\lambda\in\F_q^\times$. In either case $x\tilde tx^{-1}
\in A$, so $x$ normalizes $A$ and hence $\tilde T$. Therefore $\Cent(t,G)
\subseteq N_G(T)$. Both $T$ and $sTs^{-1}$ are abelian and contain $t$, so
both lie in $N_G(T)$. If $|T|\ge3$, the cyclic group $sTs^{-1}$ has a
generator of order at least $3$. That generator is not an involution, so it
lies in $T$, and $sTs^{-1}=T$. If $|T|=2$ then $sTs^{-1}=\{1,t\}=T$. Either
way $s\in N_G(T)$, which is a contradiction.

\emph{Step 3: counting.} Apply Theorem~\ref{thm:mackey}(a). The double
cosets inside $N_G(T)$ are $T$ and $wT$. They contribute
$\langle\pi,\pi\rangle=1$ and $\langle\pi^w,\pi\rangle=\delta$. Every other
double coset $TsT$ has $T_s=1$, so it has $|T|^2$ elements and contributes
$1$. The number of such double cosets is $(|G|-2|T|)/|T|^2$. With
$|G|=q(q-1)(q+1)$ this equals $q+2$ for $|T|=q-1$, and $q-2$ for $|T|=q+1$.
\end{proof}

\begin{counterexample}[The 2000 corollary]\label{cex:pglcor}
The 2000 draft claimed that if $\pi\neq\pi^h$ for all $h\in N_G(T)\setminus
T$, then $\Ind_T^G\pi$ is reducible. For $q\ge3$ this is true, but for a
reason unrelated to the hypothesis. By Proposition~\ref{prop:pglnorm}, the
norm is at least $q-1\ge2$ for every $\pi$. For $q=2$ it is false. Here
$G\cong S_3$, the non-split torus is $T\cong A_3$, and for $\pi\neq\triv$ we
have $\delta=0$. So the norm is $1$, and $\Ind_T^G\pi$ is the irreducible
character of Example~\ref{ex:s3ind}. This agrees with
Theorem~\ref{thm:irred}, since $\PGL(2,\F_2)\cong S_3$ is CT and $A_3$ is
its Frobenius kernel. The script confirms the split formula for $q=3,5,7$
and the non-split formula for $q=2,3$.
\end{counterexample}

%======================================================================
\section{Permutation characters and supports}\label{sec:perm}

This section corrects the main theorem of the 2000 draft. That theorem was
about linear independence of the permutation characters $\triv_A^G$.

\begin{definition}\label{def:supp}
The \emph{support} of a subgroup $H\le G$ is $\supp(H)=\{\gamma\in G_*:
\gamma\cap H\neq\emptyset\}$. Let $\tau(H)\in\F_2^{G_*}$ be its indicator
vector. Note that $\supp(H)$ is the set of classes on which $\triv_H^G$ is
non-zero.
\end{definition}

\begin{example}\label{ex:a5tau}
Order the classes of $A_5$ as $(1,2A,3A,5A,5B)$. Then $\tau(V_4)=(1,1,0,0,0)$,
$\tau(C_3)=(1,0,1,0,0)$, $\tau(C_5)=(1,0,0,1,1)$ and $\tau(A_5)=(1,1,1,1,1)$.
Over $\F_2$ we have $\tau(V_4)+\tau(C_3)+\tau(C_5)=\tau(A_5)$.
\end{example}

\begin{counterexample}[Distinct supports do not imply independence]
\label{cex:supports}
The 2000 draft argued that functions ``supported on pairwise distinct
sets'' cannot satisfy a linear relation. That is false. Take three classes
$\alpha,\beta,\gamma$, with $f_1=\triv_\alpha+\triv_\beta$,
$f_2=\triv_\alpha+\triv_\gamma$ and $f_3=2\cdot\triv_\alpha+\triv_\beta+
\triv_\gamma$. The supports are distinct, and even independent over $\F_2$.
Yet $f_3=f_1+f_2$. The converse fails too. In $A_5$ the four vectors of
Example~\ref{ex:a5tau} are dependent over $\F_2$. But the four permutation
characters are linearly independent over $\C$ (Table~\ref{tab:a5perm}).
\end{counterexample}

The correct statement needs no $\F_2$-linear algebra.

\begin{proposition}\label{prop:indep}
Let $G$ be a non-abelian commutation transitive group. Let $A_1,\dots,A_m$
represent the conjugacy classes of centralizers of non-central elements.
\begin{enumerate}[label=(\alph*),itemsep=1pt]
\item The sets $D_i=\{\gamma\in\supp(A_i):\gamma\not\subseteq Z\}$ are
  non-empty and pairwise disjoint, and their union is the set of
  non-central classes.
\item The characters $\triv_{A_1}^G,\dots,\triv_{A_m}^G,\triv$ are linearly
  independent over $\C$.
\item The vectors $\tau(A_1),\dots,\tau(A_m)$ are linearly independent over
  $\F_2$.
\end{enumerate}
\end{proposition}

\begin{proof}
(a) Let $g\notin Z$. The class of $g$ meets $A_i$ if and only if some
conjugate $g'$ of $g$ lies in $A_i\setminus Z$. By Lemma~\ref{lem:cent}(a)
this means $A_i=\Cent(g',G)$, which is conjugate to $\Cent(g,G)$. This
happens for exactly one $i$.

(b) Pick $g_i\in A_i\setminus Z$. By (a), $\triv_{A_j}^G(g_i)=0$ for
$j\neq i$. By Proposition~\ref{prop:indformula}, $\triv_{A_i}^G(g_i)=
|W(A_i)|$. Suppose $c_0\triv+\sum_jc_j\triv_{A_j}^G=0$. Evaluating at
$g_i$ gives $c_i=-c_0/|W(A_i)|$. Evaluating at $1$ then gives
$c_0(1-M)=0$, where $M=\sum_j[G:N_G(A_j)]$ is the number of distinct
centralizers of non-central elements. Now $M\ge2$, because if there were
only one, $G$ would equal $Z$ together with that abelian centralizer. So
$c_0=0$, and then every $c_i=0$.

(c) Write $e_Z$ for the indicator of the central classes. Then
$\tau(A_i)=e_Z+\triv_{D_i}$. Suppose $\sum a_i\tau(A_i)=0$. Its coordinate
on $D_i$ is $a_i$, so every $a_i$ is zero.
\end{proof}

\begin{table}[ht]
\centering
\caption{Permutation characters of $A_5$ and their decompositions. The
irreducible characters $\chi_3,\chi_3',\chi_4,\chi_5$ are those of
Table~\ref{tab:a5}.}\label{tab:a5perm}
\begin{tabular}{lccccc l}
\toprule
& $1$ & $2A$ & $3A$ & $5A$ & $5B$ & decomposition\\
\midrule
$\triv_{V_4}^G$ & 15 & 3 & 0 & 0 & 0 & $\triv+\chi_4+2\chi_5$\\
$\triv_{C_3}^G$ & 20 & 0 & 2 & 0 & 0 & $\triv+2\chi_4+\chi_5+\chi_3+\chi_3'$\\
$\triv_{C_5}^G$ & 12 & 0 & 0 & 2 & 2 & $\triv+\chi_5+\chi_3+\chi_3'$\\
$\triv_{G}^G$ & 1 & 1 & 1 & 1 & 1 & $\triv$\\
\bottomrule
\end{tabular}
\end{table}

The 2000 draft conjectured that the span of the vectors $\tau(\Cent(g,G))$
has dimension at most~$4$ for simple CT groups. The following proposition
settles this.

\begin{proposition}\label{prop:taudim}
Let $G$ be non-abelian and commutation transitive, with $m$ as in
Proposition~\ref{prop:indep}. Then $\dim_{\F_2}\Span\{\tau(\Cent(g,G)):g
\in G\}$ equals $m$ if $m$ is odd, and $m+1$ if $m$ is even. For a finite CT
group, abelian or not, it equals $1$ (abelian), $3$ (Frobenius, $m=2$) or $3$
($\PSL(2,\F_{2^f})$, $m=3$). In particular it is at most~$3$.
\end{proposition}

\begin{proof}
The vectors are $\tau(A_1),\dots,\tau(A_m)$ together with
$\tau(G)=e_Z+\sum_i\triv_{D_i}$, which comes from central $g$. Then
$\sum_i\tau(A_i)=m\,e_Z+\sum_i\triv_{D_i}$. If $m$ is odd, this equals
$\tau(G)$. If $m$ is even, writing $\tau(G)$ as a combination of the
$\tau(A_i)$ would force every coefficient to be $1$ (by the $D_i$
coordinates). The $e_Z$ coordinate would then be $m\equiv0$, which is a
contradiction. Now use Proposition~\ref{prop:indep}(c). The values of $m$
come from Propositions~\ref{prop:frobct} and~\ref{prop:psl2even}.
\end{proof}

The script confirms Proposition~\ref{prop:taudim} for $\PSL(2,\F_q)$ with
$q=4,8,16$. In each case there are $4$ distinct vectors, of rank~$3$.

\begin{yremark}\label{rem:gap}
The 2000 conjecture was based on GAP computations for $\PSL(2,\F_{2^k})$
with $2\le k\le5$, which suggested the bound~$4$. Our guess is that those
computations counted the distinct vectors ($4$) rather than the rank ($3$).
The original GAP code has not been available to check this.
\end{yremark}

%======================================================================
\section{A spanning theorem}\label{sec:span}

The 2000 draft concluded that ``all the representations in $G^*$ can be
obtained from linear combinations of monomial representations'', and
attributed this kind of result to Gelfand and Graev. That conclusion did not
follow from its theorem. This section proves a correct version, using
Brauer's induction theorem instead.

A group is \emph{elementary} if it is the direct product of a cyclic group
and a $p$-group, for some prime $p$. Elementary groups are nilpotent.

\begin{theorem}[Brauer; cited from {\cite[Ch.~10]{Se77}}]\label{thm:brauer}
Every virtual character of a finite group $G$ is a $\Z$-linear combination
of characters $\Ind_E^G\lambda$, where $E$ is an elementary subgroup and
$\lambda\in\widehat E$ is one-dimensional.
\end{theorem}

\begin{lemma}\label{lem:nilp}
In a CT group, every nilpotent subgroup is abelian. Every non-trivial
abelian subgroup $E$ lies in $\Cent(e,G)$ for any $1\neq e\in E$.
\end{lemma}

\begin{proof}
Let $E$ be nilpotent and non-abelian. Then $Z(E)\neq1$. Pick $1\neq z\in
Z(E)$. Then $E\subseteq\Cent(z,G)$, which is abelian, and this is a
contradiction. The second claim is clear.
\end{proof}

\begin{theorem}\label{thm:span}
Let $G\neq1$ be a finite CT group. Then $R(G)$ is the $\Z$-span of the
characters $\Ind_C^G\mu$, where $C=\Cent(g,G)$ for some $g\neq1$ and
$\mu\in\widehat C$.
\end{theorem}

\begin{proof}
Let $E$ be elementary and $\lambda\in\widehat E$. By
Lemma~\ref{lem:nilp}, $E$ is abelian. If $E\neq1$, choose $1\neq e\in E$
and put $C=\Cent(e,G)\supseteq E$. If $E=1$, let $C$ be any such
centralizer. Since $C$ is abelian, $\Ind_E^C\lambda$ is the sum of the
characters $\mu\in\widehat C$ that extend $\lambda$. By transitivity of induction,
$\Ind_E^G\lambda=\Ind_C^G(\Ind_E^C\lambda)$ is a non-negative integer
combination of characters $\Ind_C^G\mu$. Now apply
Theorem~\ref{thm:brauer}.
\end{proof}

\begin{table}[ht]
\centering
\caption{Character table of $A_5$, with $\phi=(1+\sqrt5)/2$. Class $5A$
contains the generator $a$ of Example~\ref{ex:a5ind}, and $5B$ contains
$a^2$.}\label{tab:a5}
\begin{tabular}{lccccc}
\toprule
& $1$ & $2A$ & $3A$ & $5A$ & $5B$\\
class size & 1 & 15 & 20 & 12 & 12\\
\midrule
$\triv$ & 1 & 1 & 1 & 1 & 1\\
$\chi_3$ & 3 & $-1$ & 0 & $\phi$ & $1-\phi$\\
$\chi_3'$ & 3 & $-1$ & 0 & $1-\phi$ & $\phi$\\
$\chi_4$ & 4 & 0 & 1 & $-1$ & $-1$\\
$\chi_5$ & 5 & 1 & $-1$ & 0 & 0\\
\bottomrule
\end{tabular}
\end{table}

\begin{example}\label{ex:a5span}
In $A_5$, let
\[
b=\Ind_{C_5}^G\mu_1,\quad b'=\Ind_{C_5}^G\mu_2,\quad
d=\Ind_{C_3}^G\omega,\quad f=\Ind_{V_4}^G\lambda,
\]
where $\omega$ and $\lambda$ are non-trivial characters of $C_3$ and $V_4$.
Proposition~\ref{prop:indformula} gives their values, and inner products
with Table~\ref{tab:a5} give
\[
b=\chi_4+\chi_5+\chi_3,\quad b'=\chi_4+\chi_5+\chi_3',\quad
d=\chi_4+2\chi_5+\chi_3+\chi_3',\quad f=\chi_4+\chi_5+\chi_3+\chi_3'.
\]
Hence
\[
\chi_3=f-b',\qquad\chi_3'=f-b,\qquad\chi_5=d-f,\qquad\chi_4=b+b'-d.
\]
This is Theorem~\ref{thm:span} made explicit. The script checks that each
of these four combinations has norm~$1$ and the right degree.

Permutation characters alone are not enough. All four in
Table~\ref{tab:a5perm} take equal values on $5A$ and $5B$, and so does every
linear combination of them. But $\chi_3$ does not. So $\chi_3$ is not in
their span, even over $\C$. The non-trivial characters $\mu$ are
essential.
\end{example}

\begin{counterexample}[$Q_8$ and the role of the center]\label{cex:q8}
Theorem~\ref{thm:span} fails for commutation transitive groups if we use
only centralizers of non-central elements. In $Q_8$ these are
$\langle i\rangle$, $\langle j\rangle$ and $\langle k\rangle$. Let
$\chi_i,\chi_j,\chi_k$ be the non-trivial linear characters, with $\chi_i$
trivial on $\langle i\rangle$, and let $\psi$ be the character of
degree~$2$. Inducing from $\langle i\rangle$ gives $\triv+\chi_i$,
$\chi_j+\chi_k$ and $\psi$ (twice). The same holds for $j$ and $k$ after
permuting the letters. Consider the homomorphism $R(Q_8)\to\Z/2$ that sends
a virtual character to the sum of its coefficients on the four linear
characters, modulo~$2$. It vanishes on every induced character, but sends
$\triv$ to~$1$. So $\triv$ is not in their $\Z$-span.
\end{counterexample}

%======================================================================
\section{Harmonic morphisms and the commuting graph}\label{sec:harmonic}

This section connects the note with graph theory. Conjugation maps the
commuting graph onto a smaller graph, the commuting conjugacy class graph.
We show that this map is a harmonic morphism. The proof rests on one fact:
the group $W(C)$ acts freely on $C\setminus Z$. On the non-commuting graph
the analogous map is almost never harmonic.

Throughout, $G$ is a finite non-abelian commutation transitive group. As in
Proposition~\ref{prop:indep}, $A_1,\dots,A_m$ represent the conjugacy
classes of centralizers of non-central elements, and $D_i$ is the set of
non-central classes meeting $A_i$.

\subsection{The commuting conjugacy class graph}

\begin{definition}[Herzog--Longobardi--Maj \cite{HLM09}]\label{def:ccc}
The \emph{commuting conjugacy class graph} $\Delta(G)$ has the non-central
conjugacy classes as vertices. Two distinct classes $\gamma,\gamma'$ are
adjacent if some $x\in\gamma$ commutes with some $y\in\gamma'$.
\end{definition}

\begin{proposition}\label{prop:ccc}
The connected components of $\Delta(G)$ are the complete graphs on
$D_1,\dots,D_m$. In particular, $\Delta(G)$ has exactly $m$ components, and
the dimension in Proposition~\ref{prop:taudim} is $m$ or $m+1$ according as
$m$ is odd or even.
\end{proposition}

\begin{proof}
Suppose $x\in\gamma$ and $y\in\gamma'$ commute, with $x,y\notin Z$. Then
$\Cent(x,G)=\Cent(y,G)$ by Lemma~\ref{lem:cent}(a). So $\gamma$ and
$\gamma'$ meet conjugates of the same $A_i$, and both lie in $D_i$.
Conversely, let $\gamma,\gamma'\in D_i$. Choose $x\in\gamma\cap A_i$ and
$y\in\gamma'\cap A_i$. Since $A_i$ is abelian, $x$ and $y$ commute. So $D_i$
spans a complete graph, and there are no edges between different $D_i$.
\end{proof}

\begin{example}\label{ex:cccexamples}
For $A_5$, the graph $\Delta(G)$ is $K_1\sqcup K_1\sqcup K_2$, with
components $\{2A\}$, $\{3A\}$ and $\{5A,5B\}$. For $\PSL(2,\F_q)$ with
$q=2^f\ge4$, it is $K_1\sqcup K_{(q-2)/2}\sqcup K_{q/2}$ by
Proposition~\ref{prop:psl2even}. For a Frobenius group $N\rtimes H$ as in
Theorem~\ref{thm:classif}(ii), it is $K_{(|N|-1)/|H|}\sqcup K_{|H|-1}$.
\end{example}

\subsection{The class map is harmonic}

The \emph{class map} $\varphi_G\colon\Gamma(G)\to\Delta(G)$ sends $x$ to its
conjugacy class. An edge $\{x,y\}$ of $\Gamma(G)$ joins commuting elements.
If $x$ and $y$ are conjugate, the edge is vertical. Otherwise it maps to the
edge joining their classes. So $\varphi_G$ is a morphism.

\begin{lemma}\label{lem:free}
Let $C=\Cent(g_0,G)$ with $g_0\notin Z$. Then $W(C)=N_G(C)/C$ acts freely
on $C\setminus Z$ by conjugation. The fibers of $\varphi_G$ on $C\setminus
Z$ are exactly the $W(C)$-orbits. So every fiber has $|W(C)|$ elements, and
$C\setminus Z$ meets $(|C|-|Z|)/|W(C)|$ classes.
\end{lemma}

\begin{proof}
Let $n\in N_G(C)$ and $c\in C\setminus Z$ with $n^{-1}cn=c$. Then
$n\in\Cent(c,G)=C$ by Lemma~\ref{lem:cent}(a). So only the identity of
$W(C)$ fixes $c$. For $c\in C\setminus Z$, the fiber of $\varphi_G$ through
$c$, intersected with $C$, is $\Conj(c,G)\cap C$. This is the
$W(C)$-orbit of $c$ by Proposition~\ref{prop:indformula}(c).
\end{proof}

\begin{theorem}\label{thm:harmonic}
Let $G$ be a finite non-abelian commutation transitive group. Then the class
map $\varphi_G\colon\Gamma(G)\to\Delta(G)$ is harmonic. More precisely, let
$C=\Cent(g_0,G)$ with $g_0\notin Z$, put $a=|W(C)|$ and
$k=(|C|-|Z|)/a$, and let $x\in C\setminus Z$.
\begin{enumerate}[label=(\alph*),itemsep=1pt]
\item The component $C\setminus Z$ of $\Gamma(G)$, a copy of $K_{ak}$, maps
  onto the component of $\Delta(G)$ containing the class of $g_0$, a copy
  of $K_k$. Its fibers all have size $a$.
\item If $k\ge2$, then $m_{\varphi_G}(x)=a$ and $v_{\varphi_G}(x)=a-1$.
  The restriction to $C\setminus Z$ has degree $a^2$, and satisfies the
  Riemann--Hurwitz formula. If $k=1$, the component collapses to a single
  vertex.
\item Let $\Delta_i$ be the component of $\Delta(G)$ on $D_i$, and suppose
  it has an edge. Its preimage is the union of the $[G:N_G(A_i)]$
  components $xA_ix^{-1}\setminus Z$. Each edge of $\Delta_i$ has
  $[G:N_G(A_i)]\cdot|W(A_i)|^2$ edges over it.
\end{enumerate}
\end{theorem}

\begin{proof}
By Proposition~\ref{prop:graph}, the component of $x$ in $\Gamma(G)$ is
$C\setminus Z$, a complete graph. By Lemma~\ref{lem:free} its image is a set
of $k$ classes, all in one $D_i$. Every class in $D_i$ meets $C$, because
$C$ is conjugate to $A_i$. So the image is the whole component $\Delta_i$,
which is complete by Proposition~\ref{prop:ccc}. The fibers have size $a$ by
Lemma~\ref{lem:free}. Now (a) and (b) follow from
Lemma~\ref{lem:collapse}(a). The edges at $x$ all lie in $C\setminus Z$, so
the restriction determines harmonicity at $x$. For (c), the distinct
conjugates of $A_i$ number $[G:N_G(A_i)]$. Each contributes $|W(A_i)|^2$
edges over a given edge of $\Delta_i$, by (b).
\end{proof}

\begin{table}[ht]
\centering
\caption{The class map $\varphi_G$ on the components of $\Gamma(G)$. Here
$q=2^f\ge4$ for $\PSL(2,\F_q)$, and $q\ge3$ for $\GL(2,\F_q)$.}
\label{tab:harmonic}
\renewcommand{\arraystretch}{1.15}
\begin{tabular}{llccc}
\toprule
$G$ & centralizer $C$ & component & image & $|W(C)|$\\
\midrule
$N\rtimes H$ (Frobenius) & $N$ & $K_{|N|-1}$ & $K_{(|N|-1)/|H|}$ & $|H|$\\
& $H$ & $K_{|H|-1}$ & $K_{|H|-1}$ & $1$\\
$\PSL(2,\F_q)$ & $U$ & $K_{q-1}$ & $K_1$ & $q-1$\\
& $T$ & $K_{q-2}$ & $K_{(q-2)/2}$ & $2$\\
& $T'$ & $K_{q}$ & $K_{q/2}$ & $2$\\
$\GL(2,\F_q)$ & $ZU$ & $K_{(q-1)^2}$ & $K_{q-1}$ & $q-1$\\
& $T$ (split) & $K_{(q-1)(q-2)}$ & $K_{(q-1)(q-2)/2}$ & $2$\\
& $T'$ (non-split) & $K_{q(q-1)}$ & $K_{q(q-1)/2}$ & $2$\\
$Q_8$, $D_8$ & each $C$ & $K_2$ & $K_1$ & $2$\\
\bottomrule
\end{tabular}
\end{table}

The entries of Table~\ref{tab:harmonic} come from
Propositions~\ref{prop:frobct} and~\ref{prop:psl2even}, and from the
following facts about $\GL(2,\F_q)$. The normalizer of a split or non-split
torus contains the torus with index~$2$, as in Step~1 of the proof of
Proposition~\ref{prop:pglnorm}. The normalizer of $ZU$ is the group $B$ of
upper triangular matrices, because $ZU$ fixes exactly one line. Conjugation
by $\diag(\alpha,\delta)$ sends $z\left(\begin{smallmatrix}1&b\\0&1
\end{smallmatrix}\right)$ to $z\left(\begin{smallmatrix}1&\alpha b/\delta\\
0&1\end{smallmatrix}\right)$. So $W(ZU)\cong B/ZU$ has order $q-1$, and it
permutes the $q-1$ values $b\ne0$ simply transitively.

\begin{example}\label{ex:a5harm}
In $A_5$, let $C=\langle a\rangle\cong C_5$. The component $C\setminus\{1\}$
is a $K_4$. Its fibers are $\{a,a^{-1}\}\subset5A$ and $\{a^2,a^{-2}\}
\subset5B$, so it maps onto the edge $\{5A,5B\}$ exactly as in
Example~\ref{ex:harmbg}. The six subgroups of order~$5$ give six such
components over the same edge. So that edge has $6\cdot4=24$ edges over it,
as Theorem~\ref{thm:harmonic}(c) predicts. The components of order~$3$ and
the Klein four-groups collapse to the single vertices $3A$ and $2A$.
\end{example}

The script checks Theorem~\ref{thm:harmonic} vertex by vertex for $S_3$,
$A_4$, $A_5$, $\PSL(2,\F_8)$, $\GL(2,\F_3)$, $Q_8$ and $D_8$. It also checks
$m_{\varphi_G}(x)=|W(\Cent(x,G))|$, $v_{\varphi_G}(x)=|W|-1$ and the
Riemann--Hurwitz formula on each component.

\subsection{The non-commuting graph}

The complement $\Gamma^c(G)$ of $\Gamma(G)$ is the non-commuting graph of
\cite{AAM06}. By Proposition~\ref{prop:graph}, it is the complete
multipartite graph whose parts are the sets $C\setminus Z$, one for each
centralizer $C$ of a non-central element. Let $M$ be the number of these
centralizers. Then $M\ge3$, because $G$ is the union of these centralizers,
and no group is a union of two proper subgroups.

\begin{proposition}\label{prop:noncommharm}
Let $\psi_G\colon\Gamma^c(G)\to K_M$ collapse each part $C\setminus Z$ to a
vertex. Then $\psi_G$ is harmonic if and only if all centralizers of
non-central elements have the same order $s+|Z|$. In that case $\deg
\psi_G=s^2$, and the Riemann--Hurwitz formula holds.
\end{proposition}

\begin{proof}
Apply Lemma~\ref{lem:collapse}(b), using $M\ge3$.
\end{proof}

\begin{corollary}\label{cor:noncommct}
If $G$ is a finite non-abelian CT group, then $\psi_G$ is not harmonic.
\end{corollary}

\begin{proof}
By Theorem~\ref{thm:classif}, $G$ is a Frobenius group or
$\PSL(2,\F_{2^f})$. In a Frobenius group $N\rtimes H$, the group $H$ acts
freely on $N\setminus\{1\}$, so $|H|$ divides $|N|-1$, and $|H|<|N|$. In
$\PSL(2,\F_q)$ the centralizers have orders $q$, $q-1$ and $q+1$.
\end{proof}

\begin{example}\label{ex:octahedron}
For $Q_8$ and $D_8$ the non-commuting graph is the octahedron
$K_{2,2,2}$, and $\psi_G$ folds it onto the triangle $K_3$ with degree $4$.
Riemann--Hurwitz reads $12=4\cdot0+6\cdot2$, since the octahedron has
genus~$7$ and the triangle has genus~$1$.

For $A_5$ even the class map on $\Gamma^c(G)$ fails to be harmonic. Its
target is the graph on the four non-trivial classes in which two classes are
adjacent when they contain non-commuting elements, here $K_4$. An
involution $x$ commutes with no element of order $3$ or $5$. So there are
$20$ edges at $x$ over $\{2A,3A\}$, but only $12$ over $\{2A,5A\}$. The
script finds that harmonicity fails at all $59$ vertices.
\end{example}

The contrast is the point of this section. Commutation transitivity makes
the commuting graph a union of cliques on which $W(C)$ acts freely. That is
exactly what harmonicity of the class map needs. The non-commuting graph
sees the different sizes of the centralizers, and harmonicity is lost.

%======================================================================
\section{Summary of results and open questions}\label{sec:summary}

This section lists the results with their status, and relates them to the
2000 draft.

\begin{center}
\renewcommand{\arraystretch}{1.15}
\begin{tabular}{>{\raggedright}p{4.6cm}>{\raggedright}p{2.2cm}>{\raggedright\arraybackslash}p{7cm}}
\toprule
Result & Status & Relation to the 2000 draft\\
\midrule
Thm.~\ref{thm:classif} (classification) & cited & Restated in full,
attributions corrected\\
Prop.~\ref{prop:gl2} ($\GL_2$, $\SL_2$) & proved & Was cited without proof\\
Cex.~\ref{cex:pgl}, Prop.~\ref{prop:psl2odd} & proved & The draft's Example
was false\\
Lemma~\ref{lem:cent}, Prop.~\ref{prop:selfnorm} & proved & Corrected
versions of draft statements\\
Props.~\ref{prop:frobct}, \ref{prop:psl2even} & proved & New\\
Prop.~\ref{prop:indformula} & proved & Draft formula off by $|C|$
(Cex.~\ref{cex:formula})\\
Thm.~\ref{thm:irred}, Cor.~\ref{cor:frobchars} & proved & New\\
Prop.~\ref{prop:pglnorm} & proved & Replaces the draft's Corollary
(Cex.~\ref{cex:pglcor})\\
Prop.~\ref{prop:indep} & proved & Replaces the draft's main theorem
(Cex.~\ref{cex:supports})\\
Prop.~\ref{prop:taudim} & proved & Settles the draft's Conjecture
(dimension $3$, not $4$)\\
Thm.~\ref{thm:span} & proved & Correct form of the draft's final claim\\
Prop.~\ref{prop:ccc}, Thm.~\ref{thm:harmonic} & proved & New (graph
theory)\\
Prop.~\ref{prop:noncommharm}, Cor.~\ref{cor:noncommct} & proved & New
(graph theory)\\
Remark~\ref{rem:gap} & unverified & Possible source of the bound $4$\\
\bottomrule
\end{tabular}
\end{center}

\begin{question}\label{q:psl}
For $G=\PSL(2,\F_{2^f})$, give a uniform formula for the decomposition of
$\Ind_C^G\mu$ into irreducibles, for every centralizer $C\in\{U,T,T'\}$ and
every $\mu\in\widehat C$. Example~\ref{ex:a5span} does this for $f=2$. The
character table of $\SL(2,\F_q)$ is classical, so this is an exercise. But a
clean statement would make Theorem~\ref{thm:span} explicit for the whole
family.
\end{question}

\begin{question}\label{q:gl2}
For commutation transitive groups with non-trivial center, what is the right
analogue of Theorem~\ref{thm:span}? Counterexample~\ref{cex:q8} shows that
centralizers of non-central elements do not suffice in general. For
$\GL(2,\F_q)$ the centralizers of non-central elements are $ZU$, $T$ and
$T'$. Characters induced from them are close to the principal series and
cuspidal representations. Is $R(\GL(2,\F_q))$ spanned by these characters
together with the linear characters of $\GL(2,\F_q)$?
\end{question}

\begin{question}\label{q:equalcent}
Which finite commutation transitive groups have all centralizers of
non-central elements of the same order? By
Proposition~\ref{prop:noncommharm}, these are exactly the groups for which
the non-commuting graph folds harmonically onto a complete graph.
Corollary~\ref{cor:noncommct} excludes every non-abelian CT group, while
$Q_8$ and $D_8$ are examples.
\end{question}

\begin{question}\label{q:corry}
The group $G$ acts on $\Gamma(G)$ by conjugation, through graph
automorphisms. Is this a harmonic group action in the sense of Corry
\cite{Co10}, and what do his genus bounds say for these graphs? The quotient
by this action is a multigraph that maps onto $\Delta(G)$, so the question
refines Theorem~\ref{thm:harmonic}.
\end{question}

\begin{question}\label{q:profinite}
Wu classified the locally finite CT groups \cite{Wu98}. Is there a version
of Theorem~\ref{thm:span} for profinite CT groups and their continuous
finite-dimensional representations? Such a representation factors through a
finite quotient. But a finite quotient of a CT group need not be CT
(compare Counterexample~\ref{cex:quotient}), so the proof does not carry
over directly.
\end{question}

%======================================================================
\section{Glossary}\label{sec:glossary}

\begin{description}[itemsep=2pt,leftmargin=0pt]
\item[AC-group.] A group whose centralizers of non-central elements are all
  abelian. Same as commutation transitive.
\item[CA-group.] A group whose centralizers of non-identity elements are all
  abelian. Same as CT.
\item[Centralizer.] $\Cent(g,G)$, the elements commuting with $g$.
\item[Character ring $R(G)$.] The $\Z$-linear combinations of irreducible
  characters. Its elements are virtual characters.
\item[Commutation transitive.] Commuting is transitive on non-central
  elements (Definition~\ref{def:ct}).
\item[Commutative transitive (CT).] Commuting is transitive on non-identity
  elements (Definition~\ref{def:ct}).
\item[Commuting conjugacy class graph.] The graph $\Delta(G)$ on the
  non-central classes, with classes adjacent when they contain commuting
  elements (Definition~\ref{def:ccc}).
\item[Commuting graph.] The graph on $G\setminus Z$ whose edges join
  commuting elements (Section~\ref{sec:graph}).
\item[Complete multipartite graph.] A graph whose vertices are split into
  parts, with two vertices adjacent exactly when they lie in different
  parts.
\item[Degree.] For a harmonic morphism, the number of edges over
  any edge of the target (Theorem~\ref{thm:bn}).
\item[Elementary group.] A direct product of a cyclic group and a
  $p$-group.
\item[Fixed-point-free automorphism.] An automorphism fixing only the
  identity.
\item[Frobenius group, kernel, complement.] See Definition~\ref{def:frob}
  and Theorem~\ref{thm:frobenius}.
\item[Genus of a graph.] $|E|-|V|+1$ for a connected graph.
\item[Harmonic morphism.] A graph morphism with constant horizontal
  multiplicity at each vertex (Definition~\ref{def:harmonic}).
\item[Horizontal multiplicity.] The number $m_\varphi(x)$ of
  Definition~\ref{def:harmonic}.
\item[Induced representation.] See \eqref{eq:indchar}.
\item[Mackey's criterion.] Theorem~\ref{thm:mackey}(b).
\item[Malnormal subgroup.] A subgroup $H$ with $H\cap xHx^{-1}=1$ for all
  $x\notin H$.
\item[Monomial representation.] A representation induced from a
  one-dimensional character.
\item[Non-commuting graph.] The complement of the commuting graph.
\item[Permutation character.] $\triv_H^G$, the number of fixed points on
  $G/H$.
\item[Riemann--Hurwitz formula.] Theorem~\ref{thm:bn}(b).
\item[Split and non-split torus.] Conjugates of the diagonal subgroup and of
  $\F_{q^2}^\times$ in $\GL(2,\F_q)$, and their images in quotients.
\item[Support.] The conjugacy classes meeting a subgroup
  (Definition~\ref{def:supp}).
\item[Unipotent element.] A matrix $g\neq1$ with $(g-1)^2=0$.
\item[Vertical edge.] An edge collapsed to a vertex by a morphism.
\item[Weyl group $W(C)$.] Here, the quotient $N_G(C)/C$ for a centralizer
  $C$.
\end{description}

%======================================================================
\section*{Acknowledgments}

The 2000 draft thanked Tony Gaglione, Dennis Spellman and Jill McGowan for
helpful conversations, and those thanks stand.

Claude (Anthropic, model Claude Opus 5.5) refereed the 2000 draft. It then
drafted this corrected and expanded version. The new material includes the
proofs of
Propositions~\ref{prop:frobct}, \ref{prop:psl2even}, \ref{prop:psl2odd},
\ref{prop:indformula}, \ref{prop:pglnorm}, \ref{prop:indep} and
\ref{prop:taudim}, \ref{prop:ccc} and \ref{prop:noncommharm}, and of
Theorems~\ref{thm:irred}, \ref{thm:span} and~\ref{thm:harmonic}. The
connection with harmonic morphisms (Section~\ref{sec:harmonic}) was
suggested by Claude and adopted by the author. The new material also
includes the literature review and the computer checks in
\texttt{comm\_tr\_checks.py}. AI review is not an independent certification
of these results. Further details on the AI usage, the sources consulted
and not consulted, and the split between original and new results are in
the revision log.

\begin{thebibliography}{FGRS00}

\bibitem[AAM06]{AAM06} A.~Abdollahi, S.~Akbari, H.~R.~Maimani,
Non-commuting graph of a group, \emph{J. Algebra} \textbf{298} (2006),
468--492.

\bibitem[BN09]{BN09} M.~Baker, S.~Norine, Harmonic morphisms and
hyperelliptic graphs, \emph{Int. Math. Res. Not. IMRN} 2009, no.~15,
2914--2955.

\bibitem[BSW58]{BSW58} R.~Brauer, M.~Suzuki, G.~E.~Wall, A characterization
of the one-dimensional unimodular projective groups over finite fields,
\emph{Illinois J. Math.} \textbf{2} (1958), 718--745.

\bibitem[Co10]{Co10} S.~Corry, Genus bounds for harmonic group actions on
finite graphs, preprint (2010), \texttt{arXiv:1006.0446}.

\bibitem[DKM09]{DKM09} D.~Dol\v{z}an, I.~Klep, P.~Moravec, Finite rings in
which commutativity is transitive, preprint (2009),
\url{https://users.fmf.uni-lj.si/moravec/Papers/ctrings23jun09.pdf}.

\bibitem[DKR05]{DKR05} A.~J.~Duncan, I.~V.~Kazachkov, V.~N.~Remeslennikov,
Centraliser dimension and universal classes of groups, preprint (2005),
\texttt{arXiv:math/0502498}.

\bibitem[FGRS00]{FGRS00} B.~Fine, A.~Gaglione, G.~Rosenberger,
D.~Spellman, The commutative transitive kernel, cited in 2000 as ``to
appear in \emph{Algebra Colloquium}''.

\bibitem[FGRS15]{FGRS15} B.~Fine, A.~Gaglione, G.~Rosenberger,
D.~Spellman, On CT and CSA groups and related ideas, preprint (2015),
\texttt{arXiv:1506.02636}.

\bibitem[Fi83]{Fi83} B.~Fine, Power conjugacy and SQ-universality in
Fuchsian and Kleinian groups, in \emph{Modular Functions in Analysis and
Number Theory}, University of Pittsburgh Press, 1983, 41--55.

\bibitem[FR09]{FR09} B.~Fine, G.~Rosenberger, Reflections on commutative
transitivity, in \emph{Aspects of Infinite Groups}, World Scientific, 2009.

\bibitem[GM23]{GM23} V.~Grazian, C.~Monetta, A conjecture related to the
nilpotency of groups with isomorphic non-commuting graphs, preprint (2023),
\texttt{arXiv:2302.01770}.

\bibitem[HLM09]{HLM09} M.~Herzog, P.~Longobardi, M.~Maj, On a commuting
graph on conjugacy classes of groups, \emph{Comm. Algebra} (2009),
\texttt{doi:10.1080/00927870802502779}.

\bibitem[Is76]{Is76} I.~M.~Isaacs, \emph{Character Theory of Finite
Groups}, Academic Press, 1976.

\bibitem[JM17]{JM17} W.~D.~Joyner, C.~G.~Melles, \emph{Adventures in Graph
Theory}, Applied and Numerical Harmonic Analysis, Birkh\"auser, 2017.

\bibitem[Jo00]{Jo00} D.~Joyner, Representation theory of $A_5$ using GAP,
preprint, 2000.

\bibitem[LR86]{LR86} F.~Levin, G.~Rosenberger, On power-commutative and
commutation-transitive groups, in \emph{Proceedings of Groups---St.\
Andrews 1985}, London Math. Soc. Lecture Note Ser. \textbf{121}, Cambridge
Univ. Press, 1986.

\bibitem[MJ26]{MJ26} C.~G.~Melles, D.~Joyner, \emph{Exploring Graphs via
Harmonic Morphisms}, Springer Nature, 2026.

\bibitem[RF23]{RF23} M.~Rezaei, Z.~Foruzanfar, The commuting conjugacy
class graphs of finite groups with a given property, \emph{Georgian Math.
J.} (2023), \texttt{doi:10.1515/gmj-2023-2069}.

\bibitem[Ro95]{Ro95} J.~J.~Rotman, \emph{An Introduction to the Theory of
Groups}, 4th ed., Grad. Texts in Math. \textbf{148}, Springer, 1995.

\bibitem[Sc70]{Sc70} R.~Schmidt, Zentralisatorverb\"ande endlicher
Gruppen, \emph{Rend. Sem. Mat. Univ. Padova} \textbf{44} (1970), 97--131.

\bibitem[Se77]{Se77} J.-P.~Serre, \emph{Linear Representations of Finite
Groups}, Grad. Texts in Math. \textbf{42}, Springer, 1977.

\bibitem[Su57]{Su57} M.~Suzuki, The nonexistence of a certain type of
simple groups of odd order, \emph{Proc. Amer. Math. Soc.} \textbf{8}
(1957), 686--695.

\bibitem[Su86]{Su86} M.~Suzuki, \emph{Group Theory II}, Grundlehren der
Math. Wiss. \textbf{248}, Springer, 1986.

\bibitem[We25]{We25} L.~Weisner, Groups in which the normaliser of every
element except identity is abelian, \emph{Bull. Amer. Math. Soc.}
\textbf{31} (1925), 413--416.

\bibitem[Ur00]{Ur00} H.~Urakawa, A discrete analogue of the harmonic
morphism and Green kernel comparison theorems, \emph{Glasgow Math. J.}
\textbf{42} (2000), 319--334.

\bibitem[Wu98]{Wu98} Y.-F.~Wu, Groups in which commutativity is a
transitive relation, \emph{J. Algebra} \textbf{207} (1998), 165--181.

\end{thebibliography}

\end{document}
